\providecommand{\uMv}{\textsf{uMv}}
\section{Conditional lower bounds}\label{sec:lower}

In this section we prove lower bounds for dynamic Betti numbers that are conditional on the \OMv\ conjecture.
Section~\ref{sec:lower-round} shows that the gadget of Section~\ref{sec:equiv} answers one \OuMv\ request with $O(\ell)$ simplex updates and a single query.
Section~\ref{sec:lower-amortized} derives amortized lower bounds for $\beta_p$ for every $p\ge1$, and Section~\ref{sec:lower-persistence} extends them to persistent Betti numbers of sublevel sets.
Section~\ref{sec:zigzag} restates the bounds as a lower bound on the total time of online zigzag persistence.
Section~\ref{sec:lower-worst} proves worst-case lower bounds for data structures that support only insertions or only deletions, and it transfers worst-case lower bounds from dynamic rank.
Section~\ref{sec:lower-bdeg} shows that the bounds hold on universes of bounded degree, Section~\ref{sec:lower-r4} on $2$-complexes that are linearly embedded in~$\R^4$, Section~\ref{sec:flag} extends them to clique complexes and Vietoris--Rips complexes, and Section~\ref{sec:acyclic} to complexes with almost no homology.
Section~\ref{sec:lower-scope} discusses what the bounds do not cover.

\subsection{One request of \OuMv}\label{sec:lower-round}

Throughout this section, \emph{dynamic $\beta_p$ over $\F$} denotes Problem~\ref{prob:betti} with parameters $p$ and~$\F$, and $n$ denotes the number of simplices of the universe.
For a vector $u\in\{0,1\}^{\ell}$ we write ${\supp u=\{i\in[\ell]:u_i=1\}}$.
A matrix with no rows or no columns has rank~$0$.

\paragraph{The hypothesis.}
We use the consequence of the \OMv\ conjecture recalled in Section~\ref{sec:prelim}~\cite[Theorem~2.7]{HKNS15full}.
If the conjecture holds, then for every constant $\delta>0$ no algorithm solves \OuMv\ for matrices $M\in\{0,1\}^{\ell\times\ell}$ and $\ell$ requests with preprocessing time polynomial in~$\ell$, computation time $O(\ell^{3-\delta})$, and error probability at most~$1/3$.
Here the \emph{computation time} is the total time spent after preprocessing.

\paragraph{Dynamic rank with switches.}
Switches make dynamic rank \OMv-hard by an elementary argument, which we state first.

\begin{lemma}\label{lem:rank-omv}
Let $\F$ be a field and $\eps>0$.
If the \OMv\ conjecture holds, then no data structure for Problem~\ref{prob:rank} over $\F$ on the instances with $r=k=\ell$ and $P=[\ell]\times[\ell]$, $\ell\ge1$, has polynomial preprocessing time, amortized update time $O(N^{1/2-\eps})$, and query time $O(N^{1-\eps})$, where $N=2\ell+\ell^2$ is the instance size.
This holds even if every entry update sets an entry to $0$ or~$1$ and all updates after the first $\ell^2$ are switches.
\end{lemma}

\begin{proof}
Given an \OuMv\ instance $M\in\{0,1\}^{\ell\times\ell}$, we set $A=M$ with at most $\ell^2$ entry updates; afterwards $R=J=[\ell]$.
For a request $(u,v)$, at most $2\ell$ switches make $R=\supp u$ and $J=\supp v$, and one query returns $\rk A[R,J]$.
Since $A[R,J]$ is a $0/1$ matrix, its rank over $\F$ is positive if and only if it has a nonzero entry, that is, if and only if $u^{\top}Mv=1$.
By the definition of amortized time, after the polynomial preprocessing of the data structure, the at most $3\ell^2$ updates and the $\ell$ interleaved queries take time $O((N+3\ell^2)N^{1/2-\eps}+\ell N^{1-\eps})=O(\ell^{3-2\eps})$ in total, and reading the requests takes $O(\ell^2)$ time; the total $O(\ell^{3-2\eps}+\ell^2)=O(\ell^{3-\delta})$, with $\delta=\min\{2\eps,1\}$, contradicts the consequence of the \OMv\ conjecture stated above.
\end{proof}

Together with Theorem~\ref{thm:equiv}(ii), which maps every switch to one simplex update and every rank query to one Betti query, Lemma~\ref{lem:rank-omv} already implies the lower bound for $\beta_1$ and $\beta_2$ in Theorem~\ref{thm:omv} below.
We nevertheless give a direct proof through the following abstraction of the gadget, because the same argument also yields the worst-case bounds for partially dynamic data structures (Theorem~\ref{thm:partial}), in which switching a row off and switching a column on are insertions, and the bounds for universes of bounded degree (Theorem~\ref{thm:bdeg}).

\paragraph{Encodings.}
Our reduction uses only three properties of the gadget, which we isolate in the following definition.
We apply it to the gadget of Section~\ref{sec:equiv} and, in Section~\ref{sec:lower-bdeg}, to a variant of bounded degree.

\begin{definition}[Encoding]\label{def:encoding}
Let $\F$ be a field, let $\ell$ be a positive integer, and let $M\in\{0,1\}^{\ell\times\ell}$.
An \emph{encoding of $M$ over~$\F$} consists of a $2$-dimensional universe $\cU$ and a subcomplex $K_{R,J}\subseteq\cU$ for every pair of sets $R,J\subseteq[\ell]$, with the following properties.
\begin{enumerate}[label=(E\arabic*)]
\item\label{enc:skeleton} Every complex $K_{R,J}$ contains all vertices and all edges of~$\cU$.
\item\label{enc:rank} $\beta_2(K_{R,J};\F)=|J|-\rk_{\F}M[R,J]$ for all $R,J\subseteq[\ell]$.
\item\label{enc:switch} There are distinct triangles $D_1,\dots,D_\ell,E_1,\dots,E_\ell$ of $\cU$, none of them in $K_{[\ell],\emptyset}$, such that $K_{R,J}=K_{[\ell],\emptyset}\cup\{D_i:i\notin R\}\cup\{E_j:j\in J\}$ for all $R,J\subseteq[\ell]$.
\end{enumerate}
\end{definition}

\begin{lemma}\label{lem:gadget-encoding}
For every field $\F$, every $\ell\ge1$, and every $M\in\{0,1\}^{\ell\times\ell}$, the universe $\cU([\ell]\times[\ell])$ of Definition~\ref{def:gadget}, together with the complexes $K_{R,J}=\Kswitch{M}{R}{J}$, is an encoding of $M$ over~$\F$, in which the triangles of~\ref{enc:switch} are the fillings $D_i$ and~$E_j$.
This universe has $\Theta(\ell^2)$ simplices and can be constructed in time $O(\ell^2)$.
\end{lemma}

\begin{proof}
Property~\ref{enc:skeleton} holds by Definition~\ref{def:gadget}, and property~\ref{enc:rank} is Lemma~\ref{lem:crux} with $A=M$ and $P=[\ell]\times[\ell]$.
By Definition~\ref{def:gadget}, the complex $\Kswitch{M}{R}{J}$ consists of $\Kswitch{M}{[\ell]}{\emptyset}$ together with the fillings $D_i$ for $i\notin R$ and $E_j$ for $j\in J$, none of which lies in $\Kswitch{M}{[\ell]}{\emptyset}$; this is~\ref{enc:switch}.
The size and the construction time are stated in Section~\ref{sec:equiv}.
\end{proof}

\begin{lemma}\label{lem:round}
Let $\cU$, together with the complexes $K_{R,J}$, be an encoding of $M\in\{0,1\}^{\ell\times\ell}$ over~$\F$.
\begin{enumerate}[label=(\roman*)]
\item For all $u,v\in\{0,1\}^{\ell}$, we have $u^{\top}Mv=1$ if and only if $\beta_2(K_{\supp u,\supp v};\F)<|\supp v|$.
\item For all $R,J\subseteq[\ell]$,
\[
  \beta_2(K_{R,J})=\beta_1(K_{R,J})-\beta_1(\cU^{(1)})+n_2(K_{R,J}),
\]
where $\cU^{(1)}$ is the $1$-skeleton of~$\cU$, so that $\beta_1(\cU^{(1)})=\beta_0(\cU^{(1)})-n_0(\cU)+n_1(\cU)$.
\item For all $R,J,R',J'\subseteq[\ell]$, the $|R\triangle R'|+|J\triangle J'|$ insertions and deletions of the triangles $D_i$, $i\in R\triangle R'$, and $E_j$, $j\in J\triangle J'$, in any order, transform $K_{R,J}$ into $K_{R',J'}$, and every intermediate family of simplices is a subcomplex of~$\cU$.
If $R'\subseteq R$ and $J\subseteq J'$, the sequence consists of insertions only; if $R\subseteq R'$ and $J'\subseteq J$, it consists of deletions only.
\end{enumerate}
\end{lemma}

\begin{proof}
(i) Let $R=\supp u$ and $J=\supp v$; Figure~\ref{fig:oumv} illustrates this part for the gadget.
By property~\ref{enc:rank}, the inequality $\beta_2(K_{R,J})<|J|$ holds if and only if $\rk_{\F}M[R,J]\ge1$, that is, if and only if the submatrix $M[R,J]$ has a nonzero entry.
Since the entries of $M$ are $0$ and $1$, and $1\ne0$ in~$\F$, this happens if and only if $M_{ij}=1$ for some $i\in R$ and $j\in J$.
This is the condition $u^{\top}Mv=1$.

(ii) Let $K=K_{R,J}$.
By~\ref{enc:skeleton}, $K$ has the same vertices and edges as~$\cU$.
Hence $n_0(K)=n_0(\cU)$ and $n_1(K)=n_1(\cU)$, and $\beta_0(K)=\beta_0(\cU^{(1)})$ because $\beta_0$ depends only on vertices and edges.
Since $\dim K\le2$, we have $\chi(K)=n_0(K)-n_1(K)+n_2(K)$, and~\eqref{eq:euler} gives
\[
  \beta_1(K)=\beta_0(\cU^{(1)})+\beta_2(K)-n_0(\cU)+n_1(\cU)-n_2(K),
\]
which is the claim.

(iii) By~\ref{enc:switch}, $K_{R',J'}$ arises from $K_{R,J}$ by inserting $D_i$ for $i\in R\setminus R'$ and $E_j$ for $j\in J'\setminus J$, and by deleting $D_i$ for $i\in R'\setminus R$ and $E_j$ for $j\in J\setminus J'$.
Every intermediate family consists of all vertices and edges of $\cU$, by~\ref{enc:skeleton}, and some of its triangles.
Since $\cU$ is $2$-dimensional, such a family is a subcomplex of~$\cU$, every inserted triangle has its edges present, and no deleted triangle has a cofacet; so every update satisfies the preconditions of Problem~\ref{prob:betti}.
If $R'\subseteq R$ and $J\subseteq J'$, all these updates are insertions; if $R\subseteq R'$ and $J'\subseteq J$, all are deletions.
\end{proof}

% Figure: answering one OuMv request with switches (Lemma~\ref{lem:round}).
% Requires only tikz. Same colors as the gadget figure: rows vermillion, columns blue.
\begin{figure}[t]
\centering
\providecommand\Description[2][]{}%
\Description{A five-by-five zero-one matrix with two rows shaded orange and two columns shaded blue; their four intersection cells are darker. Beside it a legend of switch states (hollow or filled row cycle, filled or hollow tether) and the identity relating beta two to the rank of the active submatrix.}
\definecolor{omRow}{HTML}{D55E00}%
\definecolor{omCol}{HTML}{0072B2}%
\begin{tikzpicture}[x=1cm,y=1cm,line join=round,font=\footnotesize]
  \def\s{0.44}
  % M with rows i = 1..5 (top to bottom) and columns j = 1..5
  \def\M{{{0,1,0,0,1},{1,0,0,1,0},{0,0,1,0,0},{1,0,0,0,1},{0,1,1,0,0}}}
  % highlight active rows (supp u = {1,3}) and active columns (supp v = {2,4})
  \foreach \i in {1,3} \fill[omRow!13] (0,{-\i*\s}) rectangle (5*\s,{-(\i-1)*\s});
  \foreach \j in {2,4} \fill[omCol!13] ({(\j-1)*\s},0) rectangle ({\j*\s},-5*\s);
  \foreach \i in {1,3} \foreach \j in {2,4} \fill[black!14] ({(\j-1)*\s},{-\i*\s}) rectangle ({\j*\s},{-(\i-1)*\s});
  \draw[black!30,line width=0.3pt] (0,0) grid[step=\s] (5*\s,-5*\s);
  \draw[black!70,line width=0.6pt] (0,0) rectangle (5*\s,-5*\s);
  \foreach \i in {1,...,5} \foreach \j in {1,...,5} {
    \pgfmathtruncatemacro{\e}{\M[\i-1][\j-1]}
    \ifnum\e=1 \node[inner sep=0pt,text=black] at ({(\j-0.5)*\s},{-(\i-0.5)*\s}) {$1$};
    \else \node[inner sep=0pt,text=black!40] at ({(\j-0.5)*\s},{-(\i-0.5)*\s}) {$0$}; \fi
  }
  % u and v
  \foreach \i/\e in {1/1,2/0,3/1,4/0,5/0} {
    \ifnum\e=1 \node[inner sep=0pt,text=omRow] at (-0.42,{-(\i-0.5)*\s}) {$\mathbf{1}$};
    \else \node[inner sep=0pt,text=black!45] at (-0.42,{-(\i-0.5)*\s}) {$0$}; \fi}
  \foreach \j/\e in {1/0,2/1,3/0,4/1,5/0} {
    \ifnum\e=1 \node[inner sep=0pt,text=omCol] at ({(\j-0.5)*\s},0.36) {$\mathbf{1}$};
    \else \node[inner sep=0pt,text=black!45] at ({(\j-0.5)*\s},0.36) {$0$}; \fi}
  \node[anchor=east] at (-0.62,{-2.5*\s}) {$u=$};
  \node[anchor=south] at ({2.5*\s},0.58) {$v^{\top}$};
  \node at ({2.5*\s},{-5*\s-0.35}) {$M$};
  % the switch states
  \begin{scope}[shift={(3.3,0.2)}]
    \node[anchor=west,font=\small] at (0,0) {switch states};
    \path[draw=omRow,line width=1pt,fill=white] (0.1,-0.62) -- (0.46,-0.62) -- (0.28,-0.34) -- cycle;
    \node[anchor=west] at (0.6,-0.5) {row $i\in\supp u$: $C_i$ hollow (active)};
    \path[draw=omRow,line width=1pt,fill=omRow!45] (0.1,-1.12) -- (0.46,-1.12) -- (0.28,-0.84) -- cycle;
    \node[anchor=west] at (0.6,-1.0) {row $i\notin\supp u$: filling $D_i$ present};
    \path[draw=black!80,line width=0.75pt,fill=black!28] (0.1,-1.36) -- (0.46,-1.36) -- (0.28,-1.64) -- cycle;
    \node[anchor=west] at (0.6,-1.5) {column $j\in\supp v$: filling $E_j$ present (active)};
    \path[draw=black!80,line width=0.75pt,fill=white] (0.1,-1.86) -- (0.46,-1.86) -- (0.28,-2.14) -- cycle;
    \node[anchor=west] at (0.6,-2.0) {column $j\notin\supp v$: tether $\Lambda_j$ hollow};
    \node[anchor=west,font=\small] at (0,-2.65) {$\beta_2(K)=|\supp v|-\rk M[\supp u,\supp v]$, so};
    \node[anchor=west,font=\small] at (0,-3.15) {$u^{\top}Mv=1 \iff M[\supp u,\supp v]\neq0 \iff \beta_2(K)<|\supp v|$.};
  \end{scope}
\end{tikzpicture}
\caption{Answering one \OuMv\ request $(u,v)$ with the gadget $\Kswitch{M}{R}{J}$; see Lemma~\ref{lem:round}(i).
Exactly the rows in $R=\supp u$ and the columns in $J=\supp v$ are switched on, so the active part of~$M$ is the submatrix $M[R,J]$ (dark cells).
Here $M[R,J]$ contains a~$1$, so $\rk M[R,J]=1$, $\beta_2=|J|-1<|J|$, and the Boolean product $u^{\top}Mv$ is~$1$.
Moving from one request to the next takes at most $2\ell$ switches, each of which inserts or deletes a single triangle.}
\label{fig:oumv}
\end{figure}


\subsection{Amortized lower bounds}\label{sec:lower-amortized}

\begin{theorem}\label{thm:omv}\label{thm:tradeoff}
Let $\F$ be a field, let $p\in\{1,2\}$, and let $\eps>0$.
If the \OMv\ conjecture holds, then no data structure for dynamic $\beta_p$ over $\F$ on $2$-dimensional universes has polynomial preprocessing time, amortized update time $O(n^{1/2-\eps})$, and query time $O(n^{1-\eps})$.
\end{theorem}

\begin{proof}
Suppose that $\mathcal{D}$ is such a data structure.
We use it to solve \OuMv\ for matrices $M\in\{0,1\}^{\ell\times\ell}$ and $\ell$ requests.
Let $\cU$, together with the complexes $K_{R,J}$, be the encoding of $M$ over $\F$ from Lemma~\ref{lem:gadget-encoding}, and let $n=|\cU|=\Theta(\ell^2)$.

\emph{Preprocessing.}
We construct~$\cU$, compute the number $\beta_1(\cU^{(1)})$ of Lemma~\ref{lem:round}(ii), and initialize $\mathcal{D}$ with the universe~$\cU$.
Then we insert the simplices of $K_{[\ell],\emptyset}$ into $\mathcal{D}$ in order of increasing dimension, so that every insertion is valid.
We maintain sets $R,J\subseteq[\ell]$ such that $\mathcal{D}$ stores $K_{R,J}$, initially $R=[\ell]$ and $J=\emptyset$, together with the number $n_2(K_{R,J})$ of its triangles.

\emph{Requests.}
For a request $(u,v)$, let $R'=\supp u$ and $J'=\supp v$.
We perform the at most $|R\triangle R'|+|J\triangle J'|\le2\ell$ updates of Lemma~\ref{lem:round}(iii), which transform $K_{R,J}$ into $K_{R',J'}$, and we set $R=R'$ and $J=J'$.
Then we query~$\mathcal{D}$.
If $p=2$, the answer is $\beta_2(K_{R,J})$; if $p=1$, we compute $\beta_2(K_{R,J})$ from the answer by Lemma~\ref{lem:round}(ii).
We output $1$ if $\beta_2(K_{R,J})<|J|$, and $0$ otherwise.
By Lemma~\ref{lem:round}(i), the output is $u^{\top}Mv$.

\emph{Running time.}
In total, $\mathcal{D}$ processes at most $n+2\ell^2=O(\ell^2)$ updates: at most $n$ insertions that create $K_{[\ell],\emptyset}$ and at most $2\ell$ updates per request.
By the definition of amortized time (Section~\ref{sec:prelim}), after the preprocessing of~$\mathcal{D}$ these updates together with the $\ell$ interleaved queries take total time $O((2n+2\ell^2)\cdot n^{1/2-\eps}+\ell\cdot n^{1-\eps})=O(\ell^{3-2\eps})$.
Reading a request, computing $R\triangle R'$ and $J\triangle J'$, and maintaining the counter take $O(\ell)$ time per request.
The preprocessing of $\mathcal{D}$ takes time polynomial in $n$, and hence in~$\ell$.
So the preprocessing time is polynomial in~$\ell$, and the computation time is $O(\ell^{3-2\eps}+\ell^2)=O(\ell^{3-\delta})$ for $\delta=\min\{2\eps,1\}>0$.
This contradicts the consequence of the \OMv\ conjecture stated in Section~\ref{sec:lower-round}.

We note for later use that the argument runs $\mathcal{D}$ only on the universes $\cU([\ell]\times[\ell])$, $\ell\ge1$, and that it uses these universes only through properties~\ref{enc:skeleton}--\ref{enc:switch}, the bound $n=O(\ell^2)$, and the polynomial construction time.
\end{proof}

Since worst-case bounds imply amortized bounds (Section~\ref{sec:prelim}), Theorem~\ref{thm:omv} also excludes worst-case update time $O(n^{1/2-\eps})$ together with query time $O(n^{1-\eps})$.
The reduction is deterministic.
The requests of an \OuMv\ instance may depend on the answers to earlier requests, for example when the instance is produced by the reduction from \OMv\ in~\cite{HKNS15full}; but as long as all earlier answers are correct, the requests, and hence the operations that the reduction performs on~$\mathcal{D}$, are determined by the input.
Hence the theorem also applies to randomized data structures that, for every fixed sequence of operations, answer all queries correctly with probability at least~$2/3$: we apply this guarantee to the sequence that arises when all answers are correct.
It also applies to data structures whose time bounds hold only in expectation: we run $k$ independent copies, abort a copy once its running time exceeds $12$ times its expected bound, and output the answer given by more than $k/2$ copies; by Markov's inequality, each copy answers all requests correctly within this time with probability at least $2/3-1/12=7/12>1/2$, so for a large enough constant $k$ more than $k/2$ copies do so with probability at least $2/3$.
In the operation sequences of the proof, all vertices and edges of the universe are inserted during preprocessing and are never deleted, and afterwards only triangles are inserted and deleted.
So the hardness does not come from changes of the $1$-skeleton; moreover, after preprocessing the maintained complex has $\Theta(n)$ simplices, so the bound also holds with $n$ replaced by the size of the current complex.

\paragraph{Output-sensitive data structures.}
Inserting or deleting a $q$-simplex changes exactly one Betti number, namely $\beta_q$ or $\beta_{q-1}$, and changes it by exactly one.
Indeed, inserting a $q$-simplex adds a column to $\bd_q$, which changes $\rk\bd_q$ by zero or one, and it does not change $\rk\bd_{q+1}$, because no coface of the new simplex is present; by~\eqref{eq:betti-rank}, either $\beta_q$ increases by one or $\beta_{q-1}$ decreases by one, and deletions are symmetric.
So the number of changes of the Betti numbers is exactly one per update, and Theorem~\ref{thm:omv} also excludes output-sensitive data structures whose update time is $O(f(n)+k\cdot g(n))$ with $f(n)+g(n)=O(n^{1/2-\eps})$, where $k$ is the number of changes, together with query time $O(n^{1-\eps})$.
This contrasts with time series, where the banana trees of Cultrera di Montesano et al.~\cite{CEHO24} absorb a local change in time $O(\log n+k)$, with $k$ the number of changes of the persistence diagram.

\paragraph{Higher Betti numbers.}
For $p\ge3$ we iterate the suspension of Section~\ref{sec:equiv-thm}.
For a universe $\cU$ and $m\ge0$, the \emph{$m$-fold suspension} is defined by $S^0\cU=\cU$ and $S^m\cU=S(S^{m-1}\cU)$, where each suspension uses two new vertices.
For a subcomplex $K\subseteq\cU$ we define $S^mK$ in the same way, with the same new vertices as for~$\cU$.
For $m\ge1$, the complex $S^m\emptyset$ is not empty; for instance, $S\emptyset$ consists of the two new vertices.

\begin{lemma}\label{lem:lift}
Let $\cU$ be a nonempty universe with $n$ simplices, and let $m\ge0$.
\begin{enumerate}[label=(\roman*)]
\item For every subcomplex $K\subseteq\cU$, the complex $S^mK$ is a subcomplex of $S^m\cU$, and $\beta_{q+m}(S^mK;\F)=\beta_q(K;\F)$ for every field $\F$ and every $q\ge1$.
\item $\dim S^m\cU=\dim\cU+m$, $|S^m\cU|<3^m(n+1)$, and $|S^m\emptyset|<3^m$.
\item If $K'$ arises from a subcomplex $K\subseteq\cU$ by one simplex insertion (deletion), then some sequence of $3^m$ simplex insertions (deletions) transforms $S^mK$ into $S^mK'$, and every intermediate family is a subcomplex of $S^m\cU$.
\item For fixed $m$ and every $q\ge1$, dynamic $\beta_q$ over $\F$ on $\cU$ reduces to dynamic $\beta_{q+m}$ over $\F$ on $S^m\cU$ with overhead $(3^m,1)$, and the reduction carries out insertions by insertions and deletions by deletions.
\end{enumerate}
\end{lemma}

\begin{proof}
We prove (i)--(iii) by induction on~$m$; the case $m=0$ is trivial.
Let $m\ge1$, let $L=S^{m-1}K$ and $V=S^{m-1}\cU$, so that $S^mK=SL$ and $S^m\cU=SV$.

(i) By induction $L\subseteq V$, and the suspension preserves inclusion, so $SL\subseteq SV$.
By Lemma~\ref{lem:suspension} with $p=q+m-1\ge1$, and by induction, $\beta_{q+m}(SL)=\beta_{q+m-1}(L)=\beta_q(K)$.

(ii) Since $V\supseteq\cU$ is not empty, $\dim SV=\dim V+1$.
The suspension of a complex $X$ has $|X|+2+2|X|=3|X|+2$ simplices.
By induction $|V|\le3^{m-1}(n+1)-1$, so $|SV|\le3^m(n+1)-1$.
The same computation, starting from $|\emptyset|=0$, gives $|S^m\emptyset|\le3^m-1$.

(iii) By induction, a sequence of $3^{m-1}$ updates of the same kind transforms $L$ into $S^{m-1}K'$ through subcomplexes of~$V$.
By Lemma~\ref{lem:suspension}, each of these updates is carried out by three updates of the same kind on the suspension, through subcomplexes of~$SV$.

(iv) The universe $S^m\cU$ can be computed in polynomial time, and $|S^m\cU|=O(n)$ by~(ii).
During preprocessing we insert the $3^m-1$ simplices of $S^m\emptyset$ in order of increasing dimension.
Afterwards the data structure for $S^m\cU$ stores $S^mK$.
Every update of $K$ is carried out by the $3^m$ updates of~(iii), and a query for $\beta_q(K)$ is answered by a query for $\beta_{q+m}(S^mK)$, by~(i).
\end{proof}

\begin{corollary}\label{cor:allp}
Let $\F$ be a field, let $p\ge1$ be an integer, and let $\eps>0$.
If the \OMv\ conjecture holds, then no data structure for dynamic $\beta_p$ over $\F$ on universes of dimension $\max\{p,2\}$ has polynomial preprocessing time, amortized update time $O(n^{1/2-\eps})$, and query time $O(n^{1-\eps})$.
\end{corollary}

\begin{proof}
For $p\le2$ this is Theorem~\ref{thm:omv}.
Let $p\ge3$, let $m=p-2$, and suppose that $\mathcal{D}$ is such a data structure.
For $\cU=\cU([\ell]\times[\ell])$, the universe $S^m\cU$ has dimension~$p$ and $O(n)$ simplices by Lemma~\ref{lem:lift}(ii).
By Lemma~\ref{lem:lift}(iv) and the paragraph on reductions in Section~\ref{sec:prelim}, running $\mathcal{D}$ on $S^m\cU$ yields a data structure for dynamic $\beta_2$ over $\F$ on the universes $\cU([\ell]\times[\ell])$ with polynomial preprocessing time, amortized update time $O(n^{1/2-\eps})$, and query time $O(n^{1-\eps})$.
By the proof of Theorem~\ref{thm:omv}, no such data structure exists.
\end{proof}

The dimension $\max\{p,2\}$ in Corollary~\ref{cor:allp} cannot be lowered.
If $\dim\cU<p$, then $\beta_p(K)=0$ for every subcomplex $K\subseteq\cU$.
If $\dim\cU\le1$, then $\beta_1(K)=n_1(K)-n_0(K)+\beta_0(K)$ by~\eqref{eq:euler}, and $\beta_0(K)$ is the number of connected components of the graph~$K$, which a dynamic connectivity structure maintains in $O(\log^2n)$ amortized time per update~\cite{HDT01}.
The same structure maintains $\beta_0$ of every complex (Section~\ref{sec:easy}).
So the lower bounds concern exactly the Betti numbers $\beta_p$ with $p\ge1$ on universes of dimension at least $\max\{p,2\}$.

\subsection{Persistent Betti numbers}\label{sec:lower-persistence}

We now consider sublevel-set filtrations of functions whose values change.
A function $f\colon\cU\to\R$ on a universe is \emph{monotone} if $f(\tau)\le f(\sigma)$ whenever $\tau\subseteq\sigma$.
For such $f$ and $t\in\R$, the \emph{sublevel set} ${\cU_t=\{\sigma\in\cU:f(\sigma)\le t\}}$ is a subcomplex of~$\cU$, because the faces of a simplex $\sigma\in\cU_t$ have values at most $f(\sigma)\le t$.

\begin{problem}[Sublevel-set persistence]\label{prob:sublevel}
Fix a field $\F$ and an integer $p\ge0$.
The instance is a universe $\cU$ with $n$ simplices and a monotone function $f_0\colon\cU\to\R$.
The data structure maintains a monotone function $f\colon\cU\to\R$, initially $f_0$, under the operations
\begin{itemize}
\item $\Set(\sigma,x)$, for $\sigma\in\cU$ and a rational number $x$, given in binary, such that the function obtained from $f$ by changing its value at $\sigma$ to $x$ is monotone, which makes this change;
\item $\Query(s,t)$, for real numbers $s\le t$, which returns the \emph{persistent Betti number} $\beta_p^{s,t}(f)$, the rank of the linear map $H_p(\cU_s;\F)\to H_p(\cU_t;\F)$ induced by the inclusion $\cU_s\subseteq\cU_t$.
\end{itemize}
\end{problem}

We call $\Set$ a \emph{value update}.
By the structure theorem for persistence modules~\cite{ZC05}, the vector spaces $H_p(\cU_t;\F)$, $t\in\R$, together with the maps induced by inclusions, decompose into intervals, which form the degree-$p$ \emph{barcode} of~$f$.
The persistent Betti number $\beta_p^{s,t}(f)$ is the number of intervals of the barcode that contain both $s$ and~$t$~\cite{ELZ02,EH10}.

\begin{lemma}\label{lem:sublevel}
Let $\F$ be a field and let $p\ge0$.
Dynamic $\beta_p$ over $\F$ on a universe $\cU$ reduces to Problem~\ref{prob:sublevel} over $\F$ on the same universe with overhead $(1,1)$.
In the reduction, $f$ takes values in $\{0,1\}$, and every query is $\Query(0,0)$.
\end{lemma}

\begin{proof}
For a subcomplex $K\subseteq\cU$, let $f_K(\sigma)=0$ if $\sigma\in K$ and $f_K(\sigma)=1$ otherwise.
The function $f_K$ is monotone: if $\tau\subseteq\sigma$ and $f_K(\sigma)=0$, then $\sigma\in K$, so $\tau\in K$ and $f_K(\tau)=0$; if $f_K(\sigma)=1$, then $f_K(\tau)\le f_K(\sigma)$ trivially.
Its sublevel set at~$0$ is~$K$.
We let $f_0=f_{\emptyset}$, the constant function~$1$, and we maintain the invariant $f=f_K$.
An update $\Insert(\sigma)$ produces the subcomplex $K\cup\{\sigma\}$, so $\Set(\sigma,0)$ is a valid value update that restores the invariant.
Similarly, $\Delete(\sigma)$ is carried out by $\Set(\sigma,1)$.
A query is answered by $\Query(0,0)$, which returns the rank of the identity map of $H_p(K;\F)$, that is, $\beta_p(K;\F)$.
\end{proof}

\begin{corollary}\label{cor:persistence}
Let $\F$ be a field, let $p\ge1$, and let $\eps>0$.
If the \OMv\ conjecture holds, then no data structure for Problem~\ref{prob:sublevel} over $\F$ on universes of dimension $\max\{p,2\}$ has polynomial preprocessing time, amortized update time $O(n^{1/2-\eps})$, and query time $O(n^{1-\eps})$.
This holds even when $f$ takes values in $\{0,1\}$ and every query is $\Query(0,0)$.
\end{corollary}

\begin{proof}
By Lemma~\ref{lem:sublevel} and the paragraph on reductions in Section~\ref{sec:prelim}, such a data structure yields a data structure for dynamic $\beta_p$ over $\F$ on the same universes with polynomial preprocessing time, amortized update time $O(n^{1/2-\eps})$, and query time $O(n^{1-\eps})$.
The reduction uses only values in $\{0,1\}$ and only queries $\Query(0,0)$, so a data structure that supports only these suffices.
This contradicts Corollary~\ref{cor:allp}.
\end{proof}

Corollary~\ref{cor:persistence} applies to every data structure that maintains the barcode under value updates and counts, within the stated time bounds, the intervals that contain two given numbers.
We ask for this count rather than for the barcode itself, because the barcode can have $\Theta(n)$ intervals and listing them takes linear time.
Explicit maintenance of the barcode is excluded as well.
In the reduction, $f$ takes only the values $0$ and~$1$, so the barcode of $H_p$ consists of $\beta_p(K)-\rho$ copies of $[0,1)$, $\rho$ copies of $[0,\infty)$, and $\beta_p(\cU)-\rho$ copies of $[1,\infty)$, where $K=\cU_0$ and $\rho$ is the rank of the map $H_p(K)\to H_p(\cU)$.
A value update changes $K$ by one simplex, which changes $\beta_p(K)$ by at most one and, since $H_p$ of the larger complex is spanned by the image of $H_p$ of the smaller one and at most one further class, also changes $\rho$ by at most one.
So every value update changes $O(1)$ intervals, and Corollary~\ref{cor:persistence} applies to every data structure that maintains the barcode explicitly, even an output-sensitive one.
Corollary~\ref{cor:persistence} also holds for injective functions, so data structures that assume distinct values gain nothing.
Fix an injective map $\mathrm{id}\colon\cU\to[n]$, let $d=\dim\cU$, and let $g(\sigma)=(\dim\sigma+\mathrm{id}(\sigma)/(n+1))/(2d+2)\in[0,\frac12)$.
In the proof of Lemma~\ref{lem:sublevel}, replace $f_K$ by the function that is $g(\sigma)$ on $\sigma\in K$ and $1+g(\sigma)$ on $\sigma\notin K$, and $\Query(0,0)$ by $\Query(\frac12,\frac12)$.
This function is injective, it is monotone because $g$ increases strictly from a simplex to each of its cofaces, its sublevel set at $\frac12$ is~$K$, and a simplex update of $K$ is still a single value update.
In the operation sequences obtained from the proof of Theorem~\ref{thm:omv} through Lemma~\ref{lem:sublevel}, all vertices and edges have value~$0$ after preprocessing, and only the values of triangles change.


\subsection{Online zigzag persistence}\label{sec:zigzag}

A sequence of simplex insertions and deletions is the same as a simplex-wise \emph{zigzag filtration}
\[
  \emptyset=K_0\leftrightarrow K_1\leftrightarrow\cdots\leftrightarrow K_L,
\]
in which every $K_t$ arises from $K_{t-1}$ by inserting or deleting one simplex.
Zigzag persistence~\cite{CdS10} assigns to such a filtration a barcode: a multiset of intervals of indices.
The number of intervals of the barcode of $H_p$ that contain the index $t$ is $\beta_p(K_t)$.
Algorithms for zigzag persistence go back to Carlsson, de Silva, and Morozov~\cite{CdSM09}; Maria and Oudot~\cite{MO15} update a zigzag barcode under one simplex insertion or removal at a time, in cubic total time in the worst case.
Offline, zigzag persistence can be computed in matrix multiplication time~\cite{MMS11}, and in practice by a reduction to standard persistence~\cite{DH22}; for graphs, near-linear time suffices~\cite{DH21z,DHP23}.
We call an algorithm \emph{online} if it receives the steps one at a time and, after step $t$, has reported all intervals that begin at~$t$ or end at~$t-1$.
Such an algorithm knows how many intervals of the barcode of $H_p$ contain~$t$, so it maintains $\beta_p(K_t)$ with $O(1)$ extra work per step.

\begin{corollary}[Online zigzag persistence]\label{cor:zigzag}
Assume the \OMv\ conjecture, and let $\F$ be a field and $\eps>0$.
No online algorithm computes the barcode of $H_1$ over~$\F$ of zigzag filtrations of length~$L$ of subcomplexes of a $2$-dimensional universe in total time $O(L^{3/2-\eps})$, even with polynomial preprocessing of the universe.
In particular, the amortized time per step cannot be $O(L^{1/2-\eps})$.
\end{corollary}

\begin{proof}
The reduction in the proof of Theorem~\ref{thm:omv} produces, for an \OuMv\ instance of dimension~$\ell$, a universe with $n=\Theta(\ell^2)$ simplices and a sequence of $L=\Theta(\ell^2)$ simplex insertions and deletions, interleaved with $\ell$ queries of $\beta_1$.
The steps of each round are determined by the request $(u,v)$ of that round, and the answer to the round is needed before the next request arrives; an online zigzag algorithm provides it in time.
An online zigzag algorithm answers each query in $O(1)$ time from its count of intervals containing the current index.
If its total time were $O(L^{3/2-\eps})=O(\ell^{3-2\eps})$, then \OuMv\ would be solvable in time $O(\ell^{3-2\eps}+\ell^2)=O(\ell^{3-\delta})$, with $\delta=\min\{2\eps,1\}$, after polynomial preprocessing, contradicting the \OMv\ conjecture.
\end{proof}

By Theorem~\ref{thm:acyclic}, the same bound holds for filtrations in which every complex at which the reduction reads off an answer is connected and has $\beta_1=\beta_2\le1$.
Offline, the whole barcode can be computed in time $O(L^{\omega})$~\cite{MMS11}, which exceeds $L^{3/2}$ because $\omega\ge2$, so Corollary~\ref{cor:zigzag} does not separate the online from the offline setting.
What it rules out is an online algorithm with amortized time $L^{1/2-\eps}$ per step; for graphs, in contrast, the number of intervals containing the current index is maintained online in polylogarithmic amortized time per step by dynamic connectivity (Proposition~\ref{prop:beta0}).

\subsection{Worst-case lower bounds}\label{sec:lower-worst}

\paragraph{Partially dynamic data structures.}
\emph{Incremental} dynamic $\beta_p$ over $\F$ is the variant of Problem~\ref{prob:betti} in which the instance also contains a subcomplex $K_0\subseteq\cU$, the data structure starts with $K=K_0$, and the only updates are insertions.
\emph{Decremental} dynamic $\beta_p$ over $\F$ is defined in the same way, with deletions instead of insertions.
In both variants the data structure may preprocess $K_0$ together with~$\cU$.

\begin{theorem}\label{thm:partial}
Let $\F$ be a field, let $p\ge1$, and let $\eps>0$.
If the \OMv\ conjecture holds, then no data structure for incremental dynamic $\beta_p$ over~$\F$, and no data structure for decremental dynamic $\beta_p$ over~$\F$, on universes of dimension $\max\{p,2\}$ has polynomial preprocessing time, worst-case update time $O(n^{1/2-\eps})$, and query time $O(n^{1-\eps})$.
\end{theorem}

\begin{proof}
First let $p\in\{1,2\}$.
We use the rollback technique~\cite{AV14,HKNS15}.
Suppose that $\mathcal{D}$ is an incremental data structure with the stated bounds.
Let $M\in\{0,1\}^{\ell\times\ell}$ be the matrix of an instance of \OuMv, and let $\cU$, together with the complexes $K_{R,J}$, be the encoding of $M$ over $\F$ from Lemma~\ref{lem:gadget-encoding}.

\emph{Rollback.}
We run $\mathcal{D}$ under a simulator that records, for every write to a memory cell, the address and the previous content of the cell in a log.
Replaying a log backwards restores every cell written since the log was started.
Afterwards the memory of $\mathcal{D}$, and hence its state, is the same as when the log was started.
The simulation and the log slow $\mathcal{D}$ down by a constant factor.
An operation writes at most as many cells as it takes steps, so an update and its rollback take worst-case time $O(n^{1/2-\eps})$, and a query and its rollback take worst-case time $O(n^{1-\eps})$.

\emph{Preprocessing.}
We construct $\cU$, compute~$\beta_1(\cU^{(1)})$, and initialize $\mathcal{D}$ with the universe $\cU$ and the initial complex $K_0=K_{[\ell],\emptyset}$.
This takes time polynomial in~$\ell$.

\emph{Requests.}
For a request $(u,v)$, let $R'=\supp u$ and $J'=\supp v$.
Since $R'\subseteq[\ell]$ and $\emptyset\subseteq J'$, Lemma~\ref{lem:round}(iii) gives a sequence of at most $2\ell$ triangle insertions that transforms $K_{[\ell],\emptyset}$ into $K_{R',J'}$.
We start a new log, perform these insertions, and query~$\mathcal{D}$.
If $p=1$, we compute $\beta_2(K_{R',J'})$ from the answer by Lemma~\ref{lem:round}(ii), where $n_2(K_{R',J'})$ is $n_2(K_0)$ plus the number of inserted triangles.
We output $1$ if $\beta_2(K_{R',J'})<|J'|$ and $0$ otherwise, which is $u^{\top}Mv$ by Lemma~\ref{lem:round}(i).
Finally we roll back the log.
Then $\mathcal{D}$ is again in the state it had after preprocessing, and it stores $K_{[\ell],\emptyset}$.
Between two rollbacks, $\mathcal{D}$ receives a sequence of insertions that starts from~$K_0$, so its guarantees on correctness and running time apply.

\emph{Running time.}
A request takes time $O(\ell\cdot n^{1/2-\eps}+n^{1-\eps}+\ell)=O(\ell^{2-2\eps}+\ell)$.
So the computation time is $O(\ell^{3-2\eps}+\ell^2)=O(\ell^{3-\delta})$ for $\delta=\min\{2\eps,1\}>0$, and the preprocessing time is polynomial in~$\ell$.
This contradicts the consequence of the \OMv\ conjecture stated in Section~\ref{sec:lower-round}.

For a decremental data structure we use $K_0=K_{\emptyset,[\ell]}$ instead.
Since $\emptyset\subseteq R'$ and $J'\subseteq[\ell]$, Lemma~\ref{lem:round}(iii) now gives at most $2\ell$ triangle deletions, and the rest of the argument is unchanged.
As for Theorem~\ref{thm:omv}, both arguments use only properties~\ref{enc:skeleton}--\ref{enc:switch}, the bound $n=O(\ell^2)$, and the polynomial construction time of the encodings.

Now let $p\ge3$ and $m=p-2$.
The reduction of Lemma~\ref{lem:lift}(iv) maps an initial complex $K_0$ to $S^mK_0$, insertions to insertions, and deletions to deletions.
By Lemma~\ref{lem:lift}(ii), $S^m\cU([\ell]\times[\ell])$ has dimension $p$ and $O(n)$ simplices.
Hence a data structure as in the theorem for $\beta_p$ yields an incremental, respectively decremental, data structure for $\beta_2$ on the universes $\cU([\ell]\times[\ell])$ with polynomial preprocessing time, worst-case update time $O(n^{1/2-\eps})$, and query time $O(n^{1-\eps})$, which we have just excluded.
\end{proof}

The rollback argument needs worst-case bounds.
An amortized data structure may perform an expensive operation, such as a rebuild, whose cost is charged to many earlier updates.
After a rollback, the next request can trigger the same expensive operation again, so the total time is no longer bounded by the number of updates times the amortized update time.
Theorem~\ref{thm:partial} also applies to randomized data structures that answer each query correctly with probability at least $1-1/n$.
Indeed, as long as all earlier answers are correct, each answer is produced by a run of $\mathcal{D}$ on a sequence of operations that is determined by the input, namely the preprocessing followed by the operations of one request.
By a union bound, all $\ell$ answers are correct with probability at least $1-\ell/n\ge2/3$ for all sufficiently large~$\ell$, since $n=\Theta(\ell^2)$.

\paragraph{Fully dynamic data structures.}
For fully dynamic data structures with worst-case bounds, Theorem~\ref{thm:omv} gives the lower bound discussed after its proof.
Stronger worst-case lower bounds follow from lower bounds for dynamic rank.
Recall that an entry flip replaces an entry $A_{ij}$ of a $0/1$ matrix by $1-A_{ij}$.

\begin{theorem}\label{thm:worstcase}
Let $\F$ be a field, let $p\ge1$ be an integer, and let $\alpha>0$.
Suppose that for every $\eps>0$ no data structure maintains $\rk_{\F}A$ for a matrix $A\in\{0,1\}^{\ell\times\ell}$, initially zero, under entry flips, with polynomial preprocessing time and worst-case time $O(\ell^{\alpha-\eps})$ per flip, where the time of a flip includes reporting the rank after the flip.
Then for every $\eps>0$ no data structure for dynamic $\beta_p$ over $\F$ on universes of dimension $\max\{p,2\}$ has polynomial preprocessing time, worst-case update time $O(n^{\alpha/2-\eps})$, and query time $O(n^{\alpha/2-\eps})$.
\end{theorem}

\begin{proof}
Suppose that $\mathcal{D}$ is such a data structure.
First let $p\in\{1,2\}$.
Maintaining $\rk_{\F}A$ under entry flips is the special case of Problem~\ref{prob:rank} with $r=k=\ell$, $P=[\ell]\times[\ell]$, $0/1$ entries, and no switches; its instance size is $\Theta(\ell^2)$.
By Theorem~\ref{thm:equiv}(ii), this problem reduces to dynamic $\beta_p$ over $\F$ on the $2$-dimensional universe $\cU([\ell]\times[\ell])$, which has $n=\Theta(\ell^2)$ simplices, with constant overhead.
The reduction makes $O(n)$ insertions during preprocessing, which take time $O(n\cdot n^{\alpha/2-\eps})$, polynomial in~$\ell$.
After a flip we report the rank with one query.
So $\mathcal{D}$ yields a data structure for the rank problem with polynomial preprocessing time and worst-case time $O(n^{\alpha/2-\eps})=O(\ell^{\alpha-2\eps})$ per flip, including the report of the rank.
This contradicts the hypothesis with $2\eps$ in place of~$\eps$.

For $p\ge3$, Lemma~\ref{lem:lift}(iv) with $m=p-2$ turns $\mathcal{D}$ into a data structure for dynamic $\beta_2$ over $\F$ on the universes $\cU([\ell]\times[\ell])$ with polynomial preprocessing time and worst-case update and query time $O(n^{\alpha/2-\eps})$, because $S^{m}\cU([\ell]\times[\ell])$ has dimension $p$ and $O(n)$ simplices.
The argument for $p=2$ applies to this data structure.
\end{proof}

We use the following Schur-complement identity here and in Section~\ref{sec:acyclic}.
\begin{lemma}\label{lem:bordered}
Let $Y\in\F^{m\times m}$ be invertible and $x,w\in\F^m$. Then
\[
  \rk\begin{pmatrix}Y&w\\x^{\top}&0\end{pmatrix}\;=\;m+\bigl[x^{\top}Y^{-1}w\neq0\bigr],
\]
where $[\,\cdot\,]$ is $1$ if the condition in brackets holds and $0$ otherwise.
\end{lemma}
\begin{proof}
The matrix equals \[\begin{pmatrix}I&0\\x^{\top}Y^{-1}&1\end{pmatrix}\begin{pmatrix}Y&w\\0&-x^{\top}Y^{-1}w\end{pmatrix}.\]
The first factor is invertible, and the second is block triangular with an invertible diagonal block~$Y$.
\end{proof}

The hypothesis of Theorem~\ref{thm:worstcase} follows from a hinted variant of the \OMv\ conjecture.
Van den Brand, Nanongkai, and Saranurak~\cite[Conjecture~5.12]{vdBNS19full} conjecture the hardness of the \emph{\uMv-hinted \uMv} problem~\cite[Definition~5.11]{vdBNS19full}.
In the form used in the proof of~\cite[Theorem~5.13]{vdBNS19full}, this problem maintains a product $M_1M_2M_3M_4M_5$ of five Boolean matrices of dimension at most~$s$: the matrices $M_1$, $M_3$, $M_5$ are given in advance, $M_2$ and $M_4$ start at zero and receive element updates, and a single query asks for one entry of the product.
Under the conjecture, every data structure for it with polynomial preprocessing time needs worst-case time $s^{\alpha-o(1)}$ per update or per query; we call this consequence the \emph{five-matrix bound}.
Its exponent $\alpha$ is a function of the exponents of rectangular matrix multiplication.
It satisfies $\alpha\ge5/4$, with equality if $\omega=2$, and $\alpha\approx1.407$ with the current bounds on these exponents, which also give the exponent $1.407$ of the dynamic rank algorithm of~\cite{vdBNS19} (Corollary~\ref{cor:upper} below).
As for the \OMv\ conjecture, we read the conjecture as a statement about randomized algorithms with error probability at most~$1/3$; the lower bound of van den Brand et al.\ for dynamic rank, discussed after the next proof, relies on this reading as well.

\begin{corollary}[Worst-case lower bound]\label{cor:hinted}
Let $\F$ be a field, let $p\ge1$, and let $\eps>0$.
If the \uMv-hinted \uMv\ conjecture~\cite[Conjecture~5.12]{vdBNS19full} holds, then no data structure for dynamic $\beta_p$ over $\F$ on universes of dimension $\max\{p,2\}$ has polynomial preprocessing time, worst-case update time $O(n^{\alpha/2-\eps})$, and query time $O(n^{\alpha/2-\eps})$.
Here $\alpha\ge5/4$ is the exponent of the five-matrix bound; with the current bounds on rectangular matrix multiplication, $\alpha/2\approx0.703$.
\end{corollary}

\begin{proof}
By Theorem~\ref{thm:worstcase}, it suffices to rule out a data structure $\mathcal{R}$ with the following properties.
It maintains $\rk_{\F}A$ of a matrix $A\in\{0,1\}^{\ell\times\ell}$ under entry flips.
It has polynomial preprocessing time.
It takes worst-case time $O(\ell^{\alpha-\eps})$ per flip, including the report of the rank.
By running a constant number of independent copies and taking majority votes, we may assume that $\mathcal{R}$ errs with probability at most $1/1000$ on every fixed operation sequence.

Suppose that $\mathcal{R}$ exists, and consider an instance of the five-matrix product problem above.
During preprocessing we draw independent, uniformly random $0/1$ matrices $\Xi_1$, $\Xi_3$, $\Xi_5$ of the shapes of $M_1$, $M_3$, $M_5$, and we replace $M_t$ by the entrywise product $M'_t=M_t\circ\Xi_t$ for $t\in\{1,3,5\}$; we set $M'_2=M_2$ and $M'_4=M_4$.
Following~\cite[Theorem~C.1]{vdBNS19full}, we place $M'_1,\dots,M'_5$ on the block superdiagonal of a block matrix $X$ whose diagonal blocks are identity matrices and whose other blocks are zero.
Then $X$ is a $0/1$ matrix of dimension $m=O(s)$.
It is upper unitriangular, hence invertible over every field.
The block of $X^{-1}$ in the first block row and the last block column is $-P$, where $P=M'_1M'_2M'_3M'_4M'_5$ is computed in~$\F$.
For a pair $(a,b)$, call a sequence $(i_1,i_2,i_3,i_4)$ of indices a \emph{path} if $(M_1)_{ai_1}(M_2)_{i_1i_2}(M_3)_{i_2i_3}(M_4)_{i_3i_4}(M_5)_{i_4b}=1$.
Then $P_{ab}=\sum(\Xi_1)_{ai_1}(\Xi_3)_{i_2i_3}(\Xi_5)_{i_4b}$, summed over all paths, and distinct paths contribute distinct monomials in the entries of the masks.
If entry $(a,b)$ of the Boolean product is~$0$, there is no path, and $P_{ab}=0$.
Otherwise $P_{ab}$ is a nonzero multilinear polynomial of degree three in the entries of the masks, so $P_{ab}\ne0$ with probability at least $1/8$, by the argument in the proof of Lemma~\ref{lem:isolation}.

Let $a'$ be the row of $X$ that corresponds to row $a$ of the first block, and let $b'$ be the column that corresponds to column $b$ of the last block.
By Lemma~\ref{lem:bordered} with $Y=X$, $x=e_{a'}$, and $w=e_{b'}$,
\[
  \rk\begin{pmatrix}X&e_{b'}\\e_{a'}^{\top}&0\end{pmatrix}=m+\bigl[(X^{-1})_{a'b'}\neq0\bigr]=m+\bigl[P_{ab}\ne0\bigr].
\]
We run nine independent copies of this construction, each with its own masks and its own instance of~$\mathcal{R}$ on the bordered $0/1$ matrix of dimension $\ell=m+1$, whose border row and border column are initially zero.
In preprocessing we flip the entries of each matrix $X$ into place; there are $O(s^2)$ of them, so this takes polynomial time.
An element update of $M_2$ or $M_4$ is one flip in every copy.
A query for entry $(a,b)$ sets, in every copy, the single nonzero entries of the border row and the border column, reads the rank, and resets them, which takes four flips; we answer~$1$ if and only if some copy reports rank $m+1$.
If entry $(a,b)$ of the Boolean product is~$0$, this answer is wrong only if some copy of $\mathcal{R}$ errs; otherwise it is wrong only if all nine masks fail or some copy of $\mathcal{R}$ errs.
So the error probability is at most $(7/8)^9+9/1000<1/3$.
The problem is thus solved with polynomial preprocessing, worst-case time $O(\ell^{\alpha-\eps})=O(s^{\alpha-\eps})$ per update and per query, and error probability less than $1/3$; by the proof of the five-matrix bound, this contradicts the \uMv-hinted \uMv\ conjecture.
\end{proof}

Van den Brand, Nanongkai, and Saranurak also state a lower bound for dynamic rank directly~\cite[Corollary~5.15]{vdBNS19full}.
Their rank instances, however, come from a reduction through bipartite matching that substitutes random elements of a large field for the entries, so they are not $0/1$ matrices, and our gadget encodes only matrices with entries in $\{0,\pm1\}$ (Theorem~\ref{thm:equiv}).
This is why we go back to the five-matrix product, whose block matrix~$X$ has $0/1$ entries.
Over a field of characteristic zero the masks are not needed, because $P_{ab}$ for $\Xi_1,\Xi_3,\Xi_5$ all-ones counts the paths, so that reduction is deterministic; in positive characteristic the number of paths can vanish, and the masks repair this.
Corollary~\ref{cor:upper} below gives the randomized worst-case upper bound $O(n^{1.407})$.
With the current bounds on matrix multiplication, the exponents $\alpha/2\approx0.7035$ and $\alpha\approx1.407$ differ by exactly the factor two incurred by encoding a dense $\ell\times\ell$ matrix with $\Theta(\ell^2)$ simplices.

In the other direction, dynamic-rank algorithms yield dynamic-Betti algorithms.

\begin{corollary}\label{cor:upper}
Let $\F$ be a field and let $p\ge0$ be an integer.
There is a randomized (Monte Carlo) data structure for dynamic $\beta_p$ over $\F$ that, for every fixed operation sequence of polynomial length, answers all queries correctly with high probability, with preprocessing time polynomial in~$n$, worst-case update time $O(n^{1.407})$, and query time $O(1)$, where operations in $\F$, or, if $\F$ is small, in an extension field of~$\F$ of size $n^{c}$ for a large enough constant~$c$, count as unit cost.
\end{corollary}

\begin{proof}
The block matrix of Lemma~\ref{lem:betti-to-rank} has at most $n$ rows and at most $n$ columns, and a simplex update changes at most $p+2$ of its entries (Theorem~\ref{thm:equiv}(i)).
Van den Brand, Nanongkai, and Saranurak~\cite{vdBNS19} maintain the rank of an $\ell\times\ell$ matrix under entry updates in $O(\ell^{1.407})$ worst-case time per update, with a randomized (Monte Carlo) algorithm that obtains the rank from a dynamic determinant through a reduction of Sankowski~\cite{San07}; over a small field the algorithm works in an extension field; a field of size $n^{c}$, for a large enough constant~$c$, makes the error probability of every query polynomially small, so that all answers of a polynomially long operation sequence are correct with high probability.
We run it on the block matrix, padded with zero rows and columns to size $n\times n$, and keep a counter for $n_p(K)$.
Over $\mathbb{Q}$, the unit-cost assumption can be removed: the block matrix has entries in $\{0,\pm1\}$, so by Hadamard's inequality each of its nonzero minors has absolute value at most $n^{n/2}$ and hence at most $O(n\log n)$ prime divisors.
For a prime $q$ drawn uniformly from the primes with $c\log n$ bits, for a large enough constant~$c$, the rank modulo $q$ therefore equals the rank over $\mathbb{Q}$ for all of the polynomially many matrices of a fixed operation sequence, with high probability, and arithmetic modulo $q$ takes constant time on a word RAM with words of $O(\log n)$ bits.
\end{proof}

For fully dynamic $\beta_p$, $p\ge1$, the best upper bound that we know is thus the randomized $O(n^{1.407})$ worst-case bound of Corollary~\ref{cor:upper}.
Our lower bounds are $n^{1/2-o(1)}$ for amortized, and hence also for worst-case, update time from the \OMv\ conjecture, for every data structure with query time $O(n^{1-\eps})$ (Theorem~\ref{thm:omv}), and $n^{\alpha/2-o(1)}$, with $\alpha/2\ge5/8$ and $\alpha/2\approx0.703$ for the current bounds on matrix multiplication, for worst-case update or query time over every field from the \uMv-hinted \uMv\ conjecture (Corollary~\ref{cor:hinted}).
Over $\Ftwo$, Theorem~\ref{thm:equiv} shows that closing the gap is equivalent to closing the corresponding gap for dynamic rank, with time measured in the instance size $N=r+k+m$ of Problem~\ref{prob:rank}.

\subsection{Bounded degree}\label{sec:lower-bdeg}

In the universe $\cU([\ell]\times[\ell])$, every row cycle $C_i$ meets the tubes of $\ell$ sockets, so its vertices lie in $\Theta(\ell)$ simplices.
We now replace each row cycle by a chain of copies, one for every column (Figure~\ref{fig:bdeg}).
The proof of Lemma~\ref{lem:crux} carries over by a telescoping argument.
We call the triangle cycles $s^j_0,\dots,s^j_{\ell+1}$ of column $j$ in Definition~\ref{def:gadget} the \emph{running cycles} of column~$j$; consecutive running cycles are joined by a socket.

\begin{definition}[Bounded-degree gadget]\label{def:bdeg}
Let $\ell\ge1$.
The universe $\cU^{\flat}=\cU^{\flat}(\ell)$ is obtained from $\cU([\ell]\times[\ell])$ of Definition~\ref{def:gadget} by the following changes; all new vertices are distinct.
\begin{itemize}
\item Each row cycle $C_i$ is replaced by $\ell$ triangle cycles $C_i^1,\dots,C_i^{\ell}$, the \emph{copies} of~$C_i$, and $D_i$ is the filling of~$C_i^1$.
\item For every $(i,j)\in[\ell]\times[\ell]$, the tube of the socket of position $(i,j)$ is the tube to $C_i^j$ instead of~$C_i$.
\item For every $i\in[\ell]$ and $j\in[\ell-1]$, the annulus $Q_i^j=\mathrm{Cyl}(C_i^j,C_i^{j+1})$ is added, together with its vertices and edges.
\end{itemize}
For $A\in\{0,1\}^{\ell\times\ell}$ and $R,J\subseteq[\ell]$, the complex $K^{\flat}(A;R,J)\subseteq\cU^{\flat}$ consists of all vertices and edges of~$\cU^{\flat}$, the cones, the punctured cylinders, the tubes to the tethers, all annuli $Q_i^j$, the seal of the socket of $(i,j)$ if $A_{ij}=0$ and its tube to $C_i^j$ if $A_{ij}=1$, the filling $D_i$ for every $i\notin R$, and the filling $E_j$ for every $j\in J$.
\end{definition}

% Figure: the bounded-degree gadget (Definition~\ref{def:bdeg}).
% Requires only tikz. Same colors as the gadget figure: rows vermillion, columns blue.
\begin{figure}[t]
\centering
\providecommand\Description[2][]{}%
\Description{Left: three tubes from three column surfaces meeting at a single orange row triangle. Right: the same row replaced by a horizontal pipe of four orange rings joined by annuli, each column surface tubed to its own ring, one surface sealed.}
\definecolor{bdRow}{HTML}{D55E00}%
\definecolor{bdCol}{HTML}{0072B2}%
\colorlet{bdLine}{bdCol!80!black}%
\begin{tikzpicture}[x=1cm,y=1cm,line join=round,
  surf/.style={fill=bdCol!10,draw=bdLine,line width=0.55pt},
  tube/.style={double=bdCol!24,double distance=3.6pt,line width=0.5pt,draw=bdLine},
  seal/.style={fill=black!72,draw=black,line width=0.35pt},
  hole/.style={fill=white,draw=bdLine,line width=0.45pt},
  rowcyc/.style={draw=bdRow,line width=1.1pt},
  lab/.style={font=\small},
  small/.style={font=\footnotesize}]
  % entries of row i: A_{i1}=1, A_{i2}=1, A_{i3}=0, A_{i4}=1
  %%% (a) the row cycle of Definition 4.2: all tubes of row i end on C_i
  \begin{scope}
    \draw[tube] (0.25,0.36) .. controls (0.3,1.1) and (1.35,1.3) .. (1.5,1.83);
    \draw[tube] (1.15,0.36) .. controls (1.15,1.0) and (1.55,1.3) .. (1.6,1.83);
    \draw[tube] (2.95,0.36) .. controls (2.9,1.1) and (1.85,1.3) .. (1.7,1.83);
    \foreach \x/\j in {0.25/1,1.15/2,2.05/3,2.95/4} {
      \draw[surf] (\x,0.05) ellipse (0.38 and 0.3);
      \node[small] at (\x,-0.5) {$\Sigma_{\j}$};
    }
    \foreach \x in {0.25,1.15,2.95} \draw[hole] (\x,0.33) ellipse (0.09 and 0.035);
    \draw[seal] (2.05,0.33) ellipse (0.1 and 0.045);
    \draw[rowcyc,fill=white] (1.32,1.83) -- (1.88,1.83) -- (1.6,2.25) -- cycle;
    \node[lab,text=bdRow,anchor=west] at (1.95,2.1) {$C_i$};
    \node[lab] at (1.6,-1.05) {(a)};
  \end{scope}
  %%% (b) the bounded-degree gadget: a chain of copies C_i^1,...,C_i^4 joined by annuli Q_i^j
  \begin{scope}[shift={(4.6,0)}]
    \def\cx{{0.4,1.9,3.4,4.9}}
    % the chain of annuli
    \fill[bdRow!12] (0.4,1.72) rectangle (4.9,2.28);
    \draw[bdRow!70!black,line width=0.5pt] (0.4,1.72) -- (4.9,1.72) (0.4,2.28) -- (4.9,2.28);
    \foreach \x/\j in {1.15/1,2.65/2,4.15/3} \node[small,text=bdRow!80!black] at (\x,2.0) {$Q_i^{\j}$};
    % tubes from the sockets of the 1-entries to their own copy
    \foreach \x in {0.4,1.9,4.9} \draw[tube] (\x,0.36) -- (\x,1.72);
    % copies: rings across the chain
    \foreach \x/\j in {0.4/1,1.9/2,3.4/3,4.9/4} {
      \draw[rowcyc,fill=white] (\x,2.0) ellipse (0.1 and 0.28);
      \node[small,text=bdRow,anchor=south] at (\x,2.3) {$C_i^{\j}$};
    }
    % the filling D_i of the first copy (row switch)
    \fill[bdRow!45] (0.4,2.0) ellipse (0.1 and 0.28);
    \node[small,text=black!65,anchor=east] at (0.26,2.0) {$D_i$};
    \foreach \x/\j in {0.4/1,1.9/2,3.4/3,4.9/4} {
      \draw[surf] (\x,0.05) ellipse (0.38 and 0.3);
      \node[small] at (\x,-0.5) {$\Sigma_{\j}$};
    }
    \foreach \x in {0.4,1.9,4.9} \draw[hole] (\x,0.33) ellipse (0.09 and 0.035);
    \draw[seal] (3.4,0.33) ellipse (0.1 and 0.045);
    \node[lab] at (2.65,-1.05) {(b)};
  \end{scope}
\end{tikzpicture}
\caption{Bounding the degree (Definition~\ref{def:bdeg}), shown for one row~$i$ with $A_{i1}=A_{i2}=A_{i4}=1$ and $A_{i3}=0$.
(a)~In the gadget of Definition~\ref{def:gadget}, all $\ell$ sockets of row~$i$ have a tube to the row cycle~$C_i$ in the universe (only the tubes of the $1$-entries are drawn), so every vertex of~$C_i$ lies in $\Theta(\ell)$ simplices.
(b)~In the bounded-degree gadget (here $\ell=4$), $C_i$ is replaced by a chain of copies $C_i^1,\dots,C_i^{\ell}$ joined by annuli~$Q_i^j$; the socket of position $(i,j)$ is tubed to its own copy~$C_i^j$, and the row switch $D_i$ fills the first copy.
Now every vertex of a copy lies in $O(1)$ simplices and every edge of a copy in at most three triangles, and the identity $\beta_2=|J|-\rk A[R,J]$ survives by telescoping along the chain (Lemma~\ref{lem:bdeg-crux}).}
\label{fig:bdeg}
\end{figure}


\begin{lemma}\label{lem:bdeg-crux}
Let $\F$ be a field, let $\ell\ge1$, let $A\in\{0,1\}^{\ell\times\ell}$, and let $R,J\subseteq[\ell]$.
Then
\[
  \beta_2\bigl(K^{\flat}(A;R,J);\F\bigr)=|J|-\rk_{\F}A[R,J].
\]
\end{lemma}

\begin{proof}
Let $K=K^{\flat}(A;R,J)$, and let $\Sigma_j$ be the set of triangles of $K$ that belong to column~$j$, that is, its cones, punctured cylinders, seals, and tubes.

\emph{Incidences.}
Call a vertex \emph{private} to column~$j$ if it lies on a running cycle $s^j_t$ or on the tether~$\Lambda_j$, or if it is the apex of a cone of column~$j$.
As in the proof of Lemma~\ref{lem:surface}, $\Sigma_j$ is a connected triangulated surface whose boundary is the disjoint union of $\Lambda_j$ and the copies $C_i^j$ with $A_{ij}=1$.
As there, every edge of $K$ that has an endpoint private to column $j$ and does not lie on $\Lambda_j$ lies in no triangle of $K$ or in exactly two triangles of~$K$, both in~$\Sigma_j$.
Every other edge of $K$ lies on a tether or joins two vertices of copies.
An edge that joins two vertices of copies but does not lie on a copy joins $C_i^j$ and $C_i^{j+1}$ for some $i$ and~$j$, because only the annuli contain such edges, and it lies in exactly two triangles of~$K$, both in~$Q_i^j$.
Every triangle of $K$ lies in some $\Sigma_j$ or in some $Q_i^j$, or it is a filling $D_i$ with $i\notin R$ or $E_j$ with $j\in J$.

\emph{Orientations.}
We orient all triangles as in Section~\ref{sec:crux}, and we orient each annulus $Q_i^j$ like every cylinder there.
The computation in Section~\ref{sec:crux} does not depend on the cycles to which the tubes lead.
It therefore yields a coherently oriented chain $[\Sigma_j]$ with $\bd[\Sigma_j]=[\Lambda_j]+\sum_{i:\,A_{ij}=1}[C_i^j]$, and a coherently oriented chain $[Q_i^j]$ with $\bd[Q_i^j]=[C_i^j]-[C_i^{j+1}]$.
We orient the fillings so that $\bd D_i=[C_i^1]$ and $\bd E_j=[\Lambda_j]$.

\emph{Cycles.}
Let $z$ be a $2$-cycle of~$K$.
If an edge lies in exactly two triangles $t,t'$ of $K$, both in $\Sigma_j$ or both in $Q_i^j$, then $t$ and $t'$ induce opposite orientations on it, so its coefficient in $\bd z$ is, up to sign, the difference of the coefficients of $t$ and $t'$ in~$z$; hence these coefficients agree.
The surface $\Sigma_j$ is connected through such edges, by the last claim of Lemma~\ref{lem:surface}, and so is $Q_i^j$, whose six triangles form a cyclic sequence in which consecutive triangles share an edge between $C_i^j$ and~$C_i^{j+1}$.
Therefore
\[
  z=\sum_{j\in[\ell]}x_j[\Sigma_j]+\sum_{i\in[\ell]}\sum_{j\in[\ell-1]}y_{ij}[Q_i^j]+\sum_{i\notin R}a_iD_i+\sum_{j\in J}b_jE_j
\]
for some scalars $x_j,y_{ij},a_i,b_j\in\F$.
We set $a_i=0$ for $i\in R$, $b_j=0$ for $j\notin J$, and $y_{i0}=y_{i\ell}=0$ for all~$i$.
Then
\begin{align*}
  \bd z={}&\sum_{j\in[\ell]}(x_j+b_j)\,[\Lambda_j]+\sum_{i\in[\ell]}\bigl(A_{i1}x_1+y_{i1}+a_i\bigr)\,[C_i^1]\\
  &+\sum_{i\in[\ell]}\sum_{j=2}^{\ell}\bigl(A_{ij}x_j+y_{ij}-y_{i,j-1}\bigr)\,[C_i^j].
\end{align*}
The copies and the tethers have pairwise disjoint edge sets, so $\bd z=0$ holds if and only if all coefficients vanish.

\emph{Solving.}
The coefficients of the tethers vanish if and only if $x_j=0$ for $j\notin J$ and $b_j=-x_j$ for $j\in J$.
Fix a row~$i$.
The sum of the coefficients of $C_i^1,\dots,C_i^{\ell}$ is $\sum_{j}A_{ij}x_j+a_i$, because the terms $y_{ij}$ cancel.
So if all coefficients vanish, then $\sum_{j\in J}A_{ij}x_j=0$ for $i\in R$, and $a_i=-\sum_{j}A_{ij}x_j$ for $i\notin R$.
Moreover, $y_{ij}=-a_i-\sum_{j'\le j}A_{ij'}x_{j'}$ for $j\in[\ell-1]$.
Conversely, let $x\in\F^{\ell}$ satisfy $x_j=0$ for $j\notin J$ and $\sum_{j\in J}A_{ij}x_j=0$ for $i\in R$.
If we define $a_i$, $b_j$, and $y_{ij}$ by these formulas, then the coefficients of $\Lambda_j$ and of $C_i^1,\dots,C_i^{\ell-1}$ vanish by construction, and the coefficient of $C_i^{\ell}$ vanishes because the sum of all coefficients of row $i$ is $\sum_{j}A_{ij}x_j+a_i=0$.
Hence $z\mapsto(x_j)_{j\in J}$ is a linear bijection from $Z_2(K)$ onto the kernel of $A[R,J]$.
Since $\dim K=2$, we get $\beta_2(K)=\dim Z_2(K)=|J|-\rk_{\F}A[R,J]$.
\end{proof}

\begin{theorem}\label{thm:bdeg}\label{thm:bounded-degree}
Let $\F$ be a field, let $p\in\{1,2\}$, and let $\eps>0$.
Let $\mathcal{B}$ be the class of $2$-dimensional universes in which every vertex lies in at most $37$ simplices and every edge lies in at most three triangles.
If the \OMv\ conjecture holds, then the following statements hold for universes in~$\mathcal{B}$.
\begin{enumerate}[label=(\roman*)]
\item No data structure for dynamic $\beta_p$ over $\F$ has polynomial preprocessing time, amortized update time $O(n^{1/2-\eps})$, and query time $O(n^{1-\eps})$.
\item No data structure for incremental dynamic $\beta_p$ over~$\F$, and no data structure for decremental dynamic $\beta_p$ over~$\F$, has polynomial preprocessing time, worst-case update time $O(n^{1/2-\eps})$, and query time $O(n^{1-\eps})$.
\item No data structure for Problem~\ref{prob:sublevel} over $\F$ has polynomial preprocessing time, amortized update time $O(n^{1/2-\eps})$, and query time $O(n^{1-\eps})$.
\end{enumerate}
\end{theorem}

\begin{proof}
For $\ell\ge1$ and $M\in\{0,1\}^{\ell\times\ell}$, the universe $\cU^{\flat}(\ell)$ together with the complexes $K_{R,J}=K^{\flat}(M;R,J)$ is an encoding of $M$ over~$\F$.
Indeed, \ref{enc:skeleton} holds by Definition~\ref{def:bdeg}, \ref{enc:rank} is Lemma~\ref{lem:bdeg-crux}, and \ref{enc:switch} holds with the fillings $D_i$ and $E_j$, as in $\cU([\ell]\times[\ell])$.
The universe $\cU^{\flat}(\ell)$ has $\Theta(\ell^2)$ simplices, since each position $(i,j)$ contributes one socket, one copy, and at most one annulus, each of constant size, and it can be constructed in time $O(\ell^2)$.
The proofs of Theorem~\ref{thm:omv} and of Theorem~\ref{thm:partial} for $p\in\{1,2\}$ use only these properties, so they apply to these encodings and prove (i) and~(ii) for the universes $\cU^{\flat}(\ell)$.
Part~(iii) follows from~(i) and Lemma~\ref{lem:sublevel}, which does not change the universe.
It remains to show that $\cU^{\flat}(\ell)\in\mathcal{B}$.

\emph{Triangles at a vertex.}
A vertex of a running cycle $s^j_t$ lies in at most three triangles of each of the at most two cylinders or cones that contain $s^j_t$, where a seal counts as a triangle of its cylinder.
It also lies in three triangles of each of the at most two tubes whose hole contains it.
A vertex of a copy $C_i^j$ lies in three triangles of the tube to $C_i^j$, in three triangles of each of the at most two annuli $Q_i^{j-1}$ and $Q_i^j$, and possibly in~$D_i$.
A vertex of a tether lies in three triangles of its tube and in its filling, and the apex of a cone lies in three triangles.
So every vertex lies in at most $12$ triangles.
Every edge of $\cU^{\flat}(\ell)$ lies in a triangle, so a vertex in $t$ triangles lies in at most $2t$ edges.
Hence every vertex lies in at most $1+24+12=37$ simplices.

\emph{Triangles at an edge.}
An edge of a running cycle lies in at most one triangle of each of the two cylinders or cones that contain the cycle and, if it is an edge of a hole, in one triangle of the tube of that hole.
An edge that joins two consecutive running cycles lies in at most two triangles of their cylinder and, if it is an edge of a hole, in one triangle of the tube.
An edge of a copy $C_i^j$ lies in one triangle of the tube to $C_i^j$, in one triangle of each annulus $Q_i^{j-1}$ and $Q_i^j$ that exists, and, if $j=1$, in~$D_i$; since $Q_i^0$ does not exist, these are at most three triangles.
An edge of a tether lies in one triangle of its tube and in its filling.
Every remaining edge joins the apex of a cone to its base, or a hole to the cycle at the far end of its tube, or two consecutive copies, and lies in exactly two triangles.
So every edge lies in at most three triangles.
\end{proof}

Entry flips are also local in $\cU^{\flat}(\ell)$: flipping $A_{ij}$ exchanges the seal of the socket of position $(i,j)$ with its tube to $C_i^j$, which takes a constant number of triangle updates, and Lemma~\ref{lem:bdeg-crux} holds for every matrix~$A$.
Hence the reduction in the proof of Theorem~\ref{thm:equiv}(ii), with Lemma~\ref{lem:bdeg-crux} in place of Lemma~\ref{lem:crux}, also works on the universes $\cU^{\flat}(\ell)$, and so does the proof of Theorem~\ref{thm:worstcase} for $p\in\{1,2\}$.
Since the boundary matrices of these universes are sparse, the bounds also transfer back to dynamic rank of sparse matrices.

\begin{corollary}\label{cor:bdeg-allp}
For every fixed $p\ge3$, the bounds of Theorem~\ref{thm:omv} hold for dynamic $\beta_p$ over every field~$\F$ on universes of dimension $p$ in which every vertex lies in at most $\kappa_p$ simplices, where $\kappa_p$ depends only on~$p$.
\end{corollary}

\begin{proof}
Suspension does not preserve bounded degree, because each of its two apexes lies in $\Theta(n)$ simplices.
We use a product with a sphere instead.
Let $S=\bd\Delta^{p-1}$ be the boundary of a $(p-1)$-simplex, a triangulated $(p-2)$-sphere with a constant number of simplices.
Fix a total order on the vertices of the bounded-degree universe $\cU^{\flat}$ of Theorem~\ref{thm:bdeg} and on the vertices of~$S$, and let $\cU^{\flat}\times S$ be the standard (staircase) triangulation of the product of the two ordered complexes~\cite{Mun84}.
Its simplices are the chains of the product of the two vertex orders whose projections to the factors are simplices.
Restricted to the simplices whose first projection lies in a subcomplex $K\subseteq\cU^{\flat}$, it is a triangulation of $\lVert K\rVert\times\lVert S\rVert$, which we denote $K\times S$.
Every vertex of $\cU^{\flat}\times S$ lies in a number of simplices bounded in terms of $p$ and the maximum number of simplices at a vertex of~$\cU^{\flat}$, and a simplex update of $K$ corresponds to inserting or deleting the boundedly many simplices of $K\times S$ that project onto it, in order of increasing dimension for insertions and decreasing dimension for deletions.
By the K\"unneth formula over a field~\cite[Section~3.B]{Hat02}, $H_p(K\times S)\cong\bigoplus_{i+j=p}H_i(K)\otimes H_j(S)$.
Since $\dim K=2<p$ and $H_j(S)\ne0$ only for $j\in\{0,p-2\}$, we get $\beta_p(K\times S)=\beta_2(K)\,\beta_{p-2}(S)=\beta_2(K)$.
So the reduction of Theorem~\ref{thm:bdeg} carries over with a constant factor more updates and a constant factor larger universe.
\end{proof}

\begin{corollary}\label{cor:sparse-rank}
Let $\F$ be a field and let $\eps>0$.
If the \OMv\ conjecture holds, then no data structure for Problem~\ref{prob:rank} over $\F$ has polynomial preprocessing time, amortized update time $O(N^{1/2-\eps})$, and query time $O(N^{1-\eps})$, where $N$ is the instance size.
This holds even for instances in which every row and every column contains at most three admissible positions, and for operation sequences without switches in which every entry update $\Set(i,j,x)$ has $x\in\{-1,0,1\}$.
\end{corollary}

\begin{proof}
Suppose that $\mathcal{D}$ is such a data structure.
We build a data structure for dynamic $\beta_2$ over $\F$ on the universes $\cU^{\flat}(\ell)$.
Let $B$ be the boundary matrix of $\cU^{\flat}(\ell)$ in dimension~$2$, with rows indexed by the edges and columns indexed by the triangles, for the orientations used in Lemma~\ref{lem:bdeg-crux}.
Its entries lie in $\{-1,0,1\}$.
Every column of $B$ has three nonzero entries, and every row has at most three, as shown in the proof of Theorem~\ref{thm:bdeg}.
The rank instance has $r=n_1(\cU^{\flat})$ rows, $k=n_2(\cU^{\flat})$ columns, and the admissible set $P=\supp(B)$, so its size is $N=O(n)$.
We use no switches, so all rows and columns stay active, and a rank query returns the rank of the whole matrix.
We maintain the matrix $B^K$ obtained from $B$ by replacing the columns of the triangles not in $K$ with zero columns.
Initially $K=\emptyset$ and $B^K=0$.
Inserting a triangle sets the three entries of its column to their values in~$B$, deleting a triangle sets them to~$0$, and inserting or deleting a vertex or an edge does not change~$B^K$.
Since $\cU^{\flat}$ has no $3$-simplices, Lemma~\ref{lem:betti-to-rank} gives $\beta_2(K)=n_2(K)-\rk B^K$, so a query is answered with one rank query and a counter for $n_2(K)$.
The resulting data structure has polynomial preprocessing time, amortized update time $O(n^{1/2-\eps})$, and query time $O(n^{1-\eps})$.
By the proof of Theorem~\ref{thm:bdeg}(i), no such data structure exists on the universes $\cU^{\flat}(\ell)$.
\end{proof}

Corollary~\ref{cor:sparse-rank} also follows without topology.
Replace the equation of row~$i$ by the chain $y_{i,j}-y_{i,j-1}-A_{ij}x^{(i)}_j=0$ for $j\in[\ell]$, with $y_{i,0}=0$, tie the copies of each column by $x^{(i)}_j-x^{(i-1)}_j=0$, and realize the row switch of~$i$ by the single-entry row $y_{i,\ell}=0$ and the column switch of~$j$ by the single-entry row $x^{(1)}_j=0$.
Every row and every column then has at most three nonzero entries, all in $\{0,\pm1\}$, the kernel is isomorphic to that of $A[R,J]$, and every operation changes $O(1)$ entries; together with Lemma~\ref{lem:rank-omv}, this proves Corollary~\ref{cor:sparse-rank}.
The gadget of Definition~\ref{def:bdeg} yields, in addition, the bounded-degree complexes of Theorem~\ref{thm:bdeg}.

Three admissible positions per row and column is the smallest number for which Corollary~\ref{cor:sparse-rank} can hold.
Statically, Edelsbrunner and Parsa~\cite[Theorem~4.1]{EP14}, in a special case of a sparsification of Yuster~\cite{Yus08}, reduce the rank of a $0/1$ matrix to the rank of a $0/1$ matrix with at most three nonzero entries in every row and every column, by applying their construction twice; Corollary~\ref{cor:sparse-rank} is a dynamic counterpart, and two admissible positions per column already make the problem easy.
They also remark that their complexes, in which every edge lies in at most three triangles, are already as hard as sparse rank~\cite[remark after Theorem~4.1]{EP14}, the static analogue of the class~$\mathcal{B}$ of Theorem~\ref{thm:bdeg}.

\begin{proposition}\label{rem:threshold}
Let $\F$ be a field. If every column, or every row, has at most two admissible positions, then Problem~\ref{prob:rank} over $\F$ restricted to entries in $\{-1,0,1\}$ can be solved with $O(\log^2N)$ amortized time per operation.
\end{proposition}

\begin{proof}
Let $c(\cdot)$ denote the number of connected components of a graph.
First suppose that every column has at most two admissible positions.
Let $\Gamma$ be the multigraph on the row set $[r]$ with an edge between the two nonzero rows of every active column that has two nonzero entries, a half-edge at the nonzero row of every active column that has one, and a half-edge at every inactive row.
Call an edge \emph{positive} if its two entries have opposite signs or $\F$ has characteristic~$2$, and \emph{negative} otherwise.
Then $\rk A[R,J]=\rk\bigl(A[[r],J]\;\big|\;I_{[r]\setminus R}\bigr)-(r-|R|)$, and the left kernel of the matrix on the right consists of the vectors $y\in\F^r$ with $y_i=y_{i'}$ on positive edges, $y_i=-y_{i'}$ on negative edges, and $y_i=0$ at half-edges.
Its dimension is the number of \emph{balanced} components of~$\Gamma$, those without half-edges in which every cycle has an even number of negative edges.
In the double cover $\widetilde\Gamma$, with nodes $i^{+},i^{-}$, where a positive edge $ii'$ lifts to $i^{+}i'^{+}$ and $i^{-}i'^{-}$, a negative edge to $i^{+}i'^{-}$ and $i^{-}i'^{+}$, and a half-edge at $i$ to $i^{+}i^{-}$, a component of $\Gamma$ lifts to two components if it is balanced and to one otherwise.
Hence $\rk A[R,J]=|R|-c(\widetilde\Gamma)+c(\Gamma)$.
Every operation of Problem~\ref{prob:rank} changes $O(1)$ edges of $\Gamma$ and~$\widetilde\Gamma$, so dynamic connectivity~\cite{HDT01} maintains the rank in $O(\log^2N)$ amortized time per operation; since $\rk A=\rk A^{\top}$, the same holds if every row has at most two admissible positions.
\end{proof}

This is the matrix counterpart of the contrast between $\beta_1$ of graphs, whose boundary matrices have two nonzero entries per column, and $\beta_1$ of $2$-complexes.

Over $\Ftwo$, the matrices in Corollary~\ref{cor:sparse-rank} are $0/1$ matrices.


\subsection{Universes embedded in \texorpdfstring{$\R^4$}{R4}}\label{sec:lower-r4}\label{sec:easy-r4}

Edelsbrunner and Parsa~\cite{EP14} observed that their static reduction produces $2$-complexes that embed in~$\R^4$, whereas Betti numbers of complexes in~$\R^3$ can be computed in near-linear time~\cite{DE95,DG98}.
We show that the dynamic lower bound also holds in~$\R^4$: the universes $\cU(P)$ of Definition~\ref{def:gadget} have linear embeddings in~$\R^4$, with all row cycles in a common $2$-plane and each column surface in its own \emph{page}, a closed half of a $3$-dimensional affine subspace bounded by that plane (Figure~\ref{fig:openbook}).
Section~\ref{sec:easy-s3} shows that no such lower bound can hold in~$\R^3$.
This open-book layout is that of Edelsbrunner and Parsa~\cite[Section~3]{EP14}, who place the cycles of their complex in a $2$-plane of~$\R^4$ and every punctured sphere in its own $3$-dimensional subspace, taken from a pencil of subspaces through that plane.
What is new here is that each page carries a column surface assembled from sockets, with both attachments of every socket, the tether, and its filling, at explicit rational coordinates.

A \emph{geometric simplex} in $\R^d$ is the convex hull of a set of affinely independent points, its \emph{vertices}; its \emph{faces} are the convex hulls of the subsets of its vertices, including the empty face.
A \emph{geometric simplicial complex} in $\R^d$ is a finite set $X$ of geometric simplices that contains all faces of its members and in which any two simplices intersect in a common face.
Its underlying space $\lVert X\rVert$ is the union of its simplices.
A \emph{linear embedding} of a simplicial complex $K$ in $\R^d$ is an injective map $f$ from the vertices of $K$ to $\R^d$ such that, for every $\sigma\in K$, the points $f(\sigma)$ are affinely independent, and the simplices $\operatorname{conv}f(\sigma)$, $\sigma\in K$, form a geometric simplicial complex.
For two geometric simplicial complexes $X$ and $Y$, the set $X\cap Y$ of their common simplices is a geometric simplicial complex.

\begin{lemma}[Gluing]\label{lem:glue}
Let $X$ and $Y$ be geometric simplicial complexes in $\R^d$ such that
\[\lVert X\rVert\cap\lVert Y\rVert=\lVert X\cap Y\rVert.\]
Then $X\cup Y$ is a geometric simplicial complex.
\end{lemma}

\begin{proof}
The set $X\cup Y$ contains all faces of its members.
Let $\sigma\in X$ and $\tau\in Y$.
Since $\sigma\cap\tau\subseteq\lVert X\rVert\cap\lVert Y\rVert=\lVert X\cap Y\rVert$, every point of $\sigma\cap\tau$ lies in some $\rho\in X\cap Y$.
For such~$\rho$, the set $\sigma\cap\rho$ is a common face of $\sigma$ and~$\rho$, because both lie in~$X$, and $\tau\cap\rho$ is a common face of $\tau$ and $\rho$, because both lie in~$Y$.
Both are faces of~$\rho$, so both lie in $X\cap Y$, and their intersection $\sigma\cap\tau\cap\rho$ is a common face of the two.
Hence $\sigma\cap\tau\cap\rho$ is a face of $\sigma$ and a face of~$\tau$.
So $\sigma\cap\tau$ is a union of common faces of $\sigma$ and $\tau$.
Let $V$ be the set of vertices of these faces.
Then $\sigma\cap\tau\subseteq\operatorname{conv}V$, and $\operatorname{conv}V\subseteq\sigma\cap\tau$ because $\sigma\cap\tau$ is convex and contains~$V$.
So $\sigma\cap\tau=\operatorname{conv}V$, which is a common face of $\sigma$ and~$\tau$, because $V$ consists of common vertices.
\end{proof}

\paragraph{One column.}
Fix $P\subseteq[r]\times[k]$ and a column $j\in[k]$, and let $i_1<\dots<i_d$ be the elements of $S_j$, where $d=d_j$.
We place the simplices of $\cU(P)$ that belong to column~$j$ in $\R^3$, with coordinates $(x,y,z)$.
We write the vertices of the triangle cycle $s^j_t$ as $s_{t,0},s_{t,1},s_{t,2}$, and the vertices of a triangle cycle $C$ as $C_{(0)},C_{(1)},C_{(2)}$.
Let $\delta_0=(0,1)$, $\delta_1=(1,1)$, $\delta_2=(\frac12,2)$, and $\eta_0=(0,0)$, $\eta_1=(0,1)$, $\eta_2=(1,0)$.
For $t\in\{0,\dots,d+1\}$ and $u\in\Z_3$, we place
\[
  s_{t,u}=(t,\delta_u),\qquad (\Lambda_j)_{(u)}=\bigl(\tfrac12\eta_u,\tfrac12\bigr),\qquad (C_i)_{(u)}=\bigl((i,0)+\tfrac12\eta_u,\,0\bigr)\ \ (i\in[r]),
\]
and the apexes of the cones on $s^j_0$ and on $s^j_{d+1}$ at $(-1,\frac12,\frac32)$ and $(d+2,\frac12,\frac32)$.
For $0\le t\le d$, the hole of the socket between $s^j_t$ and $s^j_{t+1}$ is the triangle cycle $h_t=(s_{t,0},s_{t,1},s_{t+1,0})$, and its vertices are the points $((t,0)+\eta_u,1)$, $u\in\Z_3$.
Let $W$ be the set of the following geometric simplices and all their faces: the six triangles of $\mathrm{Cyl}(s^j_t,s^j_{t+1})$ for $0\le t\le d$, the two cones, the tube $\mathrm{Cyl}(h_0,\Lambda_j)$, the tubes $\mathrm{Cyl}(h_t,C_{i_t})$ for $t\in[d]$, and the filling $E_j$ of~$\Lambda_j$.
Apart from the seal of $h_0$, which is not in $\cU(P)$, these are exactly the simplices of $\cU(P)$ that contain a vertex of column~$j$, together with all their faces.

\begin{lemma}\label{lem:column}
The set $W$ is a geometric simplicial complex in the half-space $\{z\ge0\}$, and $\lVert W\rVert\cap{\{z=0\}}=\bigcup_{t\in[d]}\lVert C_{i_t}\rVert$.
\end{lemma}

\begin{proof}
We assemble $W$ from three parts with Lemma~\ref{lem:glue}.

\emph{The cylinders and cones.}
Let $\Delta\subset\R^2$ be the triangle with vertices $\delta_0,\delta_1,\delta_2$, write $[\delta_u,\delta_{u+1}]$ for the segment between $\delta_u$ and $\delta_{u+1}$, and let $R_u=[0,d+1]\times[\delta_u,\delta_{u+1}]$, $u\in\Z_3$, be the lateral faces of the prism $[0,d+1]\times\Delta$.
The triangles $s_{t,u}s_{t,u+1}s_{t+1,u}$ and $s_{t,u+1}s_{t+1,u}s_{t+1,u+1}$ of the cylinders cut the rectangle $[t,t+1]\times[\delta_u,\delta_{u+1}]$ along a diagonal.
So, for fixed~$u$, the cylinder triangles in $R_u$ triangulate~$R_u$.
Two distinct faces $R_u$ meet in a segment $[0,d+1]\times\{\delta\}$, where $\delta$ is a common vertex of their edges of $\Delta$, and both triangulations subdivide this segment at the integer points.
By Lemma~\ref{lem:glue}, applied twice, the cylinder triangles and their faces form a geometric simplicial complex~$Y_0$.
The cone on $s^j_0$ consists of three facets of the tetrahedron spanned by $(-1,\frac12,\frac32)$ and $s^j_0$, so it is a geometric simplicial complex.
It lies in $\{x\le0\}$, the complex $Y_0$ lies in $\{x\ge0\}$, and both meet the plane $\{x=0\}$ exactly in the boundary of the triangle spanned by~$s^j_0$.
By Lemma~\ref{lem:glue}, and by the symmetric argument for the cone on $s^j_{d+1}$, the cylinders and the cones form a geometric simplicial complex~$Y$.
It lies in $\{1\le z\le 2\}$, and the edges of every hole $h_t$ belong to~$Y$.

\emph{The tubes.}
Let $\Delta'$ be the triangle with vertices $\eta_0,\eta_1,\eta_2$; its projection to the $x$-axis is $[0,1]$.
For $t\in[d]$ let $F_t$ be the convex hull of $h_t$ and $C_{i_t}$.
Its bottom face, spanned by $C_{i_t}$, is a copy of its top face, spanned by $h_t$, scaled by $\frac12$ and translated, with the vertices $(h_t)_{(u)}$ and $(C_{i_t})_{(u)}$ corresponding.
So $F_t$ is a frustum of a triangular pyramid.
Its lateral facets are the three trapezoids spanned by $(h_t)_{(u)},(h_t)_{(u+1)},(C_{i_t})_{(u+1)},(C_{i_t})_{(u)}$, and two of them meet exactly in a common lateral edge.
The six triangles of $\mathrm{Cyl}(h_t,C_{i_t})$ are obtained by cutting each trapezoid along its diagonal from $(h_t)_{(u+1)}$ to $(C_{i_t})_{(u)}$, so they form a geometric simplicial complex.
A point of $F_t$ at height $z\in[0,1]$ is a convex combination $zp+(1-z)q$ with $p$ in the top face and $q$ in the bottom face, so the horizontal section of $F_t$ at height $z$ is
\begin{equation}\label{eq:section}
  \tfrac{1+z}{2}\,\Delta'+\bigl(x_t(z),0\bigr),\qquad\text{where } x_t(z)=zt+(1-z)\,i_t.
\end{equation}
In the same way, the convex hull $F_0$ of $h_0$ and $\Lambda_j$ is a frustum whose section at height $z\in[\frac12,1]$ is $z\Delta'$, and $\mathrm{Cyl}(h_0,\Lambda_j)$ is a geometric simplicial complex.

\emph{The tubes meet only at shared vertices of holes.}
For $z<1$ the section~\eqref{eq:section} projects to the $x$-interval $[x_t(z),x_t(z)+\frac{1+z}{2}]$, whose length is less than one.
For $1\le t<t'\le d$ we have $t'-t\ge1$ and $i_{t'}-i_t\ge1$, so $x_{t'}(z)-x_t(z)=z(t'-t)+(1-z)(i_{t'}-i_t)\ge1$.
Hence $F_t$ and $F_{t'}$ can meet only at height $z=1$, where they meet in $\operatorname{conv}h_t\cap\operatorname{conv}h_{t'}$.
This is the point $s_{t+1,0}$ if $t'=t+1$, and empty otherwise.
Similarly, the section of $F_0$ at height $z\in[\frac12,1)$ projects into $[0,z]$, while $x_t(z)\ge1$ for $t\ge1$.
So $F_0$ meets $F_t$ only in the point $s_{1,0}$, and only if $t=1$.
These points are common vertices of the tubes involved.
By Lemma~\ref{lem:glue}, applied once for each tube, the tubes and their faces form a geometric simplicial complex~$Z$.

\emph{Gluing.}
The complex $Y$ lies in $\{z\ge1\}$, and $Z$ lies in $\{z\le1\}$.
Each tube meets the plane $\{z=1\}$ exactly in the boundary of the triangle spanned by its hole, and the edges of this boundary belong to $Y$ and to~$Z$.
So $\lVert Y\rVert\cap\lVert Z\rVert=\lVert Y\cap Z\rVert$, and $Y\cup Z$ is a geometric simplicial complex by Lemma~\ref{lem:glue}.
The filling $E_j$ is the bottom face of $F_0$ at height~$\frac12$.
It does not meet $Y$, it meets $\mathrm{Cyl}(h_0,\Lambda_j)$ exactly in $\lVert \Lambda_j\rVert$, and it does not meet $F_t$ for $t\ge1$, because $E_j$ projects into $[0,\frac12]$ while the section of $F_t$ at height $\frac12$ projects into $[1,\infty)$.
A final application of Lemma~\ref{lem:glue} shows that $W$ is a geometric simplicial complex.
All its vertices have $z\ge0$, so $W$ lies in $\{z\ge0\}$.
Only the tubes $\mathrm{Cyl}(h_t,C_{i_t})$ with $t\in[d]$ reach height $0$, and they meet $\{z=0\}$ exactly in $\lVert C_{i_t}\rVert$.
\end{proof}

% Figure: the open-book embedding of the hard universes in R^4 (Theorem~\ref{thm:r4}).
% Requires only tikz. Same colors as the gadget figure: rows vermillion, columns blue.
\begin{figure}[t]
\centering
\providecommand\Description[2][]{}%
\Description{Left: three pages fanning out from a common horizontal binding line that carries three small orange triangles. Right: the three pages drawn flat side by side, each with a blue band (column surface) near the top, blue tubes down to some of the triangles on the binding, dark seals on the band above the others, and a small gray tether.}
\definecolor{obRow}{HTML}{D55E00}%
\definecolor{obCol}{HTML}{0072B2}%
\colorlet{obLine}{obCol!80!black}%
\begin{tikzpicture}[x=1cm,y=1cm,line join=round,
  page/.style={fill=black!3,draw=black!40,line width=0.5pt},
  band/.style={fill=obCol!12,draw=obLine,line width=0.6pt,rounded corners=2pt},
  tube/.style={fill=obCol!26,draw=obLine,line width=0.45pt},
  seal/.style={fill=black!75,draw=black,line width=0.3pt},
  rowtri/.style={fill=white,draw=obRow,line width=1pt},
  teth/.style={fill=black!28,draw=black!80,line width=0.6pt},
  binding/.style={black!75,line width=1.5pt}]
  %%% (a) the book: three pages (half-spaces) with the common boundary Pi
  \begin{scope}[shift={(0,0.1)}]
    \foreach \vx/\vy in {-0.55/1.3,0.55/1.3,0/1.65}
      \path[page,fill opacity=0.8] (0,0) -- (2.2,0) -- (2.2+\vx,\vy) -- (\vx,\vy) -- cycle;
    \node[font=\small,anchor=south] at (-0.55,1.3) {$\Sigma_1$};
    \node[font=\small,anchor=south] at (1.1,1.65) {$\Sigma_2$};
    \node[font=\small,anchor=south] at (2.75,1.3) {$\Sigma_3$};
    \draw[binding] (-0.1,0) -- (2.3,0);
    \foreach \x in {0.5,1.1,1.7} \path[rowtri] (\x-0.11,-0.07) -- (\x+0.11,-0.07) -- (\x,0.11) -- cycle;
    \node[font=\small,anchor=north] at (1.1,-0.14) {$\Pi$};
    \node[font=\small] at (1.1,-0.8) {(a)};
  \end{scope}
  %%% (b) the contents of the three pages, each drawn flat; their bottom edges are the same plane Pi
  % \obpanel{x-offset}{j}{entry row 1}{entry row 2}{entry row 3}
  \def\obpanel#1#2#3#4#5{%
    \begin{scope}[shift={(#1,0)}]
      \path[page] (0,0) rectangle (2.75,2.0);
      \foreach \x/\e in {0.7/#3,1.375/#4,2.05/#5} {\ifnum\e=1 \path[tube] (\x-0.08,0.08) rectangle (\x+0.08,1.31);\fi}
      \path[tube] (0.34,0.9) rectangle (0.5,1.31);
      \path[band] (0.22,1.3) rectangle (2.53,1.58);
      \foreach \x/\e in {0.7/#3,1.375/#4,2.05/#5} {\ifnum\e=0 \path[seal] (\x-0.12,1.34) -- (\x+0.12,1.34) -- (\x,1.54) -- cycle;\fi}
      \path[teth] (0.27,0.9) -- (0.57,0.9) -- (0.42,0.72) -- cycle;
      \node[font=\small] at (2.2,1.8) {$\Sigma_{#2}$};
      \draw[binding] (-0.06,0) -- (2.81,0);
      \foreach \x/\i in {0.7/1,1.375/2,2.05/3} {
        \path[rowtri] (\x-0.12,-0.08) -- (\x+0.12,-0.08) -- (\x,0.12) -- cycle;
        \node[font=\scriptsize,text=obRow,anchor=north] at (\x,-0.12) {$C_{\i}$};
      }
    \end{scope}}
  \obpanel{3.5}{1}{1}{0}{1}
  \obpanel{6.55}{2}{1}{1}{0}
  \obpanel{9.6}{3}{0}{1}{1}
  \node[font=\small] at (7.925,-0.7) {(b)};
\end{tikzpicture}
\caption{The open-book embedding of Theorem~\ref{thm:r4}, an explicit linear version of the pencil-of-$3$-planes embedding of Edelsbrunner and Parsa~\cite{EP14}, drawn schematically for the $3\times3$ matrix~$A$ with columns $(1,0,1)$, $(1,1,0)$, and $(0,1,1)$.
(a)~The row cycles $C_1,C_2,C_3$ lie in the binding plane~$\Pi$, and each column surface $\Sigma_j$ lies in its own page, a closed half of a $3$-dimensional affine subspace of~$\R^4$ bounded by~$\Pi$; distinct pages meet only in~$\Pi$.
(b)~Each page is drawn flat, with its bottom edge representing~$\Pi$.
In the page of~$\Sigma_j$, a tube joins the hole of the socket of position $(i,j)$ to $C_i$ if $A_{ij}=1$, and a seal (black) closes the hole if $A_{ij}=0$; a short tube leads to the private tether~$\Lambda_j$ (gray).
So different surfaces meet only along the row cycles.
Here $P=[3]\times[3]$; in the universe $\cU(P)$, the hole of every socket has both its seal and its tube.}
\label{fig:openbook}
\end{figure}


\paragraph{The open book.}
For $j\in[k]$ let $\varphi_j\colon\R^3\to\R^4$ be the linear map $\varphi_j(x,y,z)=(x,y,z,jz)$, and let $\Pi=\{(x,y,0,0):x,y\in\R\}$.
Each $\varphi_j$ is injective and fixes the plane $\{z=0\}$ in the sense that $\varphi_j(x,y,0)=(x,y,0,0)$.
If $j\ne j'$ and $\varphi_j(x,y,z)=\varphi_{j'}(x',y',z')$, then $x=x'$, $y=y'$, $z=z'$, and $jz=j'z$, so $z=0$.
Hence the images of distinct maps $\varphi_j$ meet exactly in~$\Pi$.
Figure~\ref{fig:openbook} shows the resulting embedding.


\begin{theorem}\label{thm:r4}
\leavevmode
\begin{enumerate}[label=(\roman*)]
\item For all $r,k\ge1$ and every $P\subseteq[r]\times[k]$, the complex $\cU(P)$ of Definition~\ref{def:gadget} has a linear embedding in $\R^4$ with rational coordinates, and such an embedding can be computed in ${O(r+k+|P|)}$ time.
\item Assume the \OMv\ conjecture, and let $\F$ be a field, $p\in\{1,2\}$, and $\eps>0$.
No data structure with polynomial preprocessing time solves Problem~\ref{prob:betti} for $\beta_p$ over $\F$ with amortized update time $O(n^{1/2-\eps})$ and query time $O(n^{1-\eps})$, even if the universe is a $2$-dimensional geometric simplicial complex in $\R^4$ that is given with the coordinates of its vertices.
\end{enumerate}
\end{theorem}

\begin{proof}
(i)~For every column $j\in[k]$, let $W_j$ be the complex of Lemma~\ref{lem:column}, and let $W'_j$ be $W_j$ without the seal of $h_0$.
Removing a triangle that is a face of no other simplex keeps a geometric simplicial complex, so $W'_j$ is one.
Since $\varphi_j$ is linear and injective, $\varphi_j(W'_j)$ is a geometric simplicial complex in~$\R^4$.
The position of a vertex of a row cycle does not depend on the column, because $\varphi_j(x,y,0)=(x,y,0,0)$ for every~$j$.
We place the filling $D_i$ of $C_i$ in $\Pi$, for every $i\in[r]$.
The triangles $D_i$ are pairwise disjoint, because $D_i$ projects to the $x$-interval $[i,i+\frac12]$.
Let $X$ be the union of the complexes $\varphi_j(W'_j)$, $j\in[k]$, and of the triangles $D_i$, $i\in[r]$, with their faces.

We show by Lemma~\ref{lem:glue} that $X$ is a geometric simplicial complex, adding one column at a time and then the triangles~$D_i$.
By Lemma~\ref{lem:column}, the complex $\varphi_j(W'_j)$ meets $\Pi$ exactly in $\bigcup_{i\in S_j}\lVert C_i\rVert$.
The images of distinct maps $\varphi_j$ meet only in~$\Pi$.
So, for $j\ge2$, the union of the first $j-1$ columns meets $\varphi_j(W'_j)$ exactly in $\bigcup\lVert C_i\rVert$, over the rows $i\in S_j$ that also lie in $S_{j'}$ for some $j'<j$.
This set is the underlying space of the common simplices, namely the vertices and edges of these row cycles.
Likewise, the union of all columns meets the triangles $D_i$ exactly in the row cycles that occur in some column, and these consist of common simplices.
So $X$ is a geometric simplicial complex.

The vertex positions are injective: the vertices of column $j$ that are not on a row cycle have $z\ge\frac12$, so they are mapped into the image of $\varphi_j$ minus~$\Pi$, and distinct vertices of $W_j$ have distinct positions.
The simplices of $X$ are the images of the simplices of $\cU(P)$: those that contain a vertex of column $j$ come from $W'_j$, and the remaining ones are the row cycles and their fillings.
So the vertex positions form a linear embedding of $\cU(P)$.
All coordinates are rational, and each is computed in $O(1)$ time; $\cU(P)$ has $O(r+k+|P|)$ simplices.

(ii)~The proof of Theorem~\ref{thm:omv} uses only universes of the form $\cU([\ell]\times[\ell])$ (Lemma~\ref{lem:gadget-encoding}).
By~(i), the coordinates of a linear embedding of such a universe in $\R^4$ can be computed during preprocessing, in linear time, and added to the instance.
So the proof of Theorem~\ref{thm:omv} applies without change.
\end{proof}

\begin{corollary}\label{cor:r4-equiv}
Let $\F$ be a field.
Dynamic rank over $\F$ of $0/1$ matrices reduces with overhead $(7,1)$ to dynamic $\beta_1$, and to dynamic $\beta_2$, over $\F$ on $2$-complexes given with a linear embedding in~$\R^4$.
Over $\Ftwo$, these problems are equivalent to dynamic rank with constant overhead.
\end{corollary}

\begin{proof}
The reduction of Theorem~\ref{thm:equiv}(ii) uses the universe $\cU(P)$, and by Theorem~\ref{thm:r4}(i) its linear embedding can be computed in linear time and added to the instance.
Conversely, dynamic $\beta_p$ on every universe, embedded or not, reduces to dynamic rank by Theorem~\ref{thm:equiv}(i), and over $\Ftwo$ every matrix is a $0/1$ matrix.
\end{proof}
Over other fields we do not know whether the universes $\cU^{\pm}(P)$ with reversed tubes, which encode $\{0,\pm1\}$-matrices, embed in~$\R^4$.


\subsection{Clique complexes and Vietoris--Rips complexes}\label{sec:flag}

Many complexes in topological data analysis are flag complexes, Vietoris--Rips complexes for example.
The \emph{clique complex} $\Cl(G)$ of a graph $G$ has one simplex for every nonempty clique of~$G$.
The \emph{Vietoris--Rips complex} $\VR_\varrho(X,\mathrm{dist})$ of a finite metric space $(X,\mathrm{dist})$ at scale $\varrho$ is the clique complex of the graph that joins two points at distance at most~$\varrho$; for a finite point set $X\subset\R^d$ with the Euclidean metric we write $\VR_\varrho(X)$.
Such a complex is determined by its graph, so the natural updates insert or delete a vertex together with its incident edges, or insert or delete a point.
One such update may add or remove many simplices at once.
We show that the lower bounds survive in these models.

\begin{problem}[Dynamic clique-complex Betti numbers]\label{prob:clique}
Fix a field $\F$ and $p\ge0$.
The instance is a graph $G=(V,E)$ with $n$ vertices.
The data structure maintains a set $W\subseteq V$, initially empty, under insertion and deletion of single vertices, and on query returns $\beta_p(\Cl(G[W]);\F)$, where $G[W]$ is the subgraph induced by~$W$.
\end{problem}

\begin{problem}[Dynamic Vietoris--Rips Betti numbers]\label{prob:rips}
Fix a field $\F$ and an integer $p\ge0$.
The instance is a dimension $d\ge1$ and a set $X_{\cU}\subset\mathbb{Q}^d$ of $n$ points.
The data structure maintains a set $X\subseteq X_{\cU}$, initially empty, under insertion and deletion of single points, and on query returns $\beta_p(\VR_1(X);\F)$.
\end{problem}

\begin{theorem}\label{thm:flag}
Assume the \OMv\ conjecture, and let $\F$ be a field and $\eps>0$.
For each of the following problems, no data structure with polynomial preprocessing time has amortized update time $O(n^{1/2-\eps})$ and query time $O(n^{1-\eps})$.
In each problem, a query returns the first Betti number over~$\F$ of the complex named.
\begin{enumerate}[label=(\roman*)]
\item\label{it:flag-vertex} The clique complex $\Cl(G[W])$, where $G$ is a fixed graph with $n$ vertices and maximum degree bounded by an absolute constant, and an update inserts a vertex into $W$ or deletes one from it (Problem~\ref{prob:clique} with $p=1$).
\item\label{it:flag-edge} The clique complex $\Cl(H)$, where $H$ is a spanning subgraph of a fixed graph $G$ as in~\ref{it:flag-vertex}, and an update inserts an edge of $G$ into $H$ or deletes one from it.
\item\label{it:flag-metric} The Vietoris--Rips complex $\VR_1(X,\mathrm{dist})$, where $X$ is a fixed set of $n$ points, $\mathrm{dist}$ is a metric on $X$ with values in $\{1,2\}$, and an update changes the distance between one pair of points.
\item\label{it:flag-euclid} The Vietoris--Rips complex $\VR_1(X)$ of Problem~\ref{prob:rips} with $d=O(\log n)$, even if all coordinates of the points have $O(\log n)$ bits.
\end{enumerate}
\end{theorem}

The proof rests on barycentric subdivision.
For a simplicial complex $\cU$ let $G_{\cU}$ be the \emph{comparability graph} of its face poset.
Its vertices are the simplices of~$\cU$, and two simplices are adjacent when one is a proper face of the other.
For a subcomplex $K\subseteq\cU$, the induced subgraph $G_{\cU}[K]$ is the comparability graph of the face poset of~$K$.
Its cliques are exactly the chains $\sigma_0\subsetneq\dots\subsetneq\sigma_q$ of simplices of~$K$.
So $\Cl(G_{\cU}[K])$ is the order complex of the face poset of~$K$, which is the barycentric subdivision $\mathrm{sd}(K)$~\cite{Koz08}.
Since $\mathrm{sd}(K)$ is homeomorphic to~$K$,
\begin{equation}\label{eq:sd}
  \beta_p\bigl(\Cl(G_{\cU}[K])\bigr)\;=\;\beta_p(K)\qquad\text{for every subcomplex }K\subseteq\cU\text{ and every }p.
\end{equation}

\begin{proof}[Proof of Theorem~\ref{thm:flag}]
Let $\cU$ be the bounded-degree universe of Theorem~\ref{thm:bounded-degree}.
It has $n$ simplices, every vertex lies in $O(1)$ simplices, and the lower bound of Theorem~\ref{thm:bounded-degree}(i) for $\beta_1$ holds on it.
In $G_{\cU}$ a simplex is adjacent only to its faces and to the simplices that contain it, so $G_{\cU}$ has maximum degree $O(1)$.

For~\ref{it:flag-vertex} we maintain $W=K$.
Inserting or deleting a simplex $\sigma$ of $K$ becomes inserting or deleting the vertex $\sigma$ of~$W$.
By~\eqref{eq:sd} the queries agree.

For~\ref{it:flag-edge} we maintain the spanning subgraph $H$ of $G_{\cU}$ whose edges are the edges of $G_{\cU}[K]$; the simplices of $\cU$ that are not in $K$ are isolated vertices of~$H$.
When $\sigma$ is inserted into~$K$, all its proper faces lie in $K$ and none of its cofaces does, so the new edges of $H$ are the edges between $\sigma$ and its proper faces; there are at most six of them, because $\dim\cU=2$.
A deletion removes the same edges.
So every simplex update becomes at most six edge updates, and we query only after all of them have been carried out.
At that moment $\Cl(H)$ is the disjoint union of $\Cl(G_{\cU}[K])$ and some isolated vertices, so $\beta_1(\Cl(H))=\beta_1(K)$ by~\eqref{eq:sd}.

For~\ref{it:flag-metric} we take $X$ to be the vertex set of $G_{\cU}$ and set $\mathrm{dist}(x,x')=1$ if $xx'$ is an edge of the graph $H$ of~\ref{it:flag-edge}, and $\mathrm{dist}(x,x')=2$ for all other pairs of distinct points.
Every symmetric function with values in $\{1,2\}$ satisfies the triangle inequality, because $2\le1+1$, so $\mathrm{dist}$ is a metric.
By definition, $\VR_1(X,\mathrm{dist})=\Cl(H)$, and an edge update of $H$ is a single distance update.

For~\ref{it:flag-euclid} we realize $G_{\cU}$ as a unit-ball graph: we map its vertices to points so that adjacent vertices are at distance less than~$1$ and nonadjacent vertices at distance greater than~$1$.
Let $A_G$ be the adjacency matrix of $G=G_{\cU}$ and $\Delta$ its maximum degree.
All eigenvalues of $A_G$ lie in $[-\Delta,\Delta]$, so $A_G+\lambda I$ with $\lambda=\Delta+1$ is positive definite.
Write it as a Gram matrix with vectors $y_x\in\R^{n}$, one per vertex.
Then $\lVert y_x-y_{x'}\rVert^2=2\lambda-2$ if $xx'$ is an edge and $2\lambda$ otherwise.
By the Johnson--Lindenstrauss lemma~\cite{DG03}, a random linear map $\Phi\colon\R^{n}\to\R^{d}$ with $d=O(\lambda^2\log n)$ preserves all these squared distances up to a factor $1\pm1/(4\lambda)$ with high probability.
It can be checked in polynomial time; we redraw it at most $n$ times, and if all draws fail, which happens with probability $2^{-\Omega(n)}$, the reduction gives up, which only adds this probability to its error.
Edges then have squared length below $2\lambda-1$ and non-edges above it.
Set $f(x)=\Phi(y_x)/\sqrt{2\lambda-1}$.
Then $\lVert f(x)-f(x')\rVert<1$ for edges and $>1$ for non-edges.
This is the argument of Frankl and Maehara~\cite{FM88}, who bound the sphericity of a graph in terms of its least eigenvalue.
For every $W\subseteq V$, $\VR_1(f(W))=\Cl(G_{\cU}[W])$ as abstract complexes, because no pairwise distance equals~$1$.
Inserting or deleting the vertex $\sigma$ thus becomes inserting or deleting the point $f(\sigma)$.
Since $\Delta=O(1)$, we get $d=O(\log n)$.
Because the squared distances are separated from the threshold by a constant margin after scaling, rounding all coordinates to rationals with $O(\log n)$ bits does not change the graph, so the points have a polynomial-size description.
\end{proof}

A query of the form $\beta_1(\VR_\varrho(X))$ at a fixed scale $\varrho$ is a persistent Betti number of the Vietoris--Rips filtration.
Theorem~\ref{thm:flag} therefore also bounds every data structure that maintains the $H_1$-persistence of Vietoris--Rips filtrations of a dynamic point cloud, or of a finite metric space whose distances change one at a time, and can report the number of bars alive at a fixed scale.
Applied to the bounded-degree universes of Theorem~\ref{thm:bdeg} for $p=2$ and of Corollary~\ref{cor:bdeg-allp} for $p\ge3$, the same proof gives statements~\ref{it:flag-vertex}--\ref{it:flag-euclid} for $\beta_p$, for every fixed $p\ge1$, with $2^{p+1}-2$ in place of six in~\ref{it:flag-edge}; Theorem~\ref{thm:acyclic} transfers in the same way to statement~\ref{it:flag-metric}, and to statements~\ref{it:flag-vertex} and~\ref{it:flag-edge} without the degree bound, because the cone vertex of its universe lies in $\Theta(n)$ simplices.
Two questions remain open: whether the Vietoris--Rips statement holds in a fixed dimension $d\ge3$ (Section~\ref{sec:conclusion}), and whether it holds under point \emph{moves} at bounded speed rather than insertions and deletions.

\subsection{Hardness with almost no homology}\label{sec:acyclic}

The complexes built in the proof of Theorem~\ref{thm:omv} carry a lot of homology, because they contain the whole $1$-skeleton of the universe.
Even without the unused edges the homology stays large.
Let $A\in\{0,1\}^{r\times k}$, and let $K^-(A)$ be the subcomplex of $K(A;[r],[k])$ that consists of its triangles, all their faces, and the row cycles.
We call it the \emph{lean} complex.
In $K^-(A)$ all rows are active and all tethers are filled.
The surface $\Sigma_j$ together with the filling $E_j$ of its tether is a sphere with one hole for every nonzero entry of column~$j$, and the row cycles are circles, so
\begin{equation}\label{eq:euler-KA}
  \chi(K^-(A))\;=\;\sum_{j=1}^{k}\bigl(2-\nnz(A_{\ast j})\bigr)\;=\;2k-\nnz(A),
\end{equation}
where $\nnz(\cdot)$ counts nonzero entries and $A_{\ast j}$ is column~$j$ of~$A$.
Let $G_A$ denote the \emph{support graph} of $A$: the bipartite graph with a node $\rho_i$ for every row, a node $\gamma_j$ for every column, and an edge $\rho_i\gamma_j$ for every nonzero entry $A_{ij}$.
Write $c(G_A)$ for its number of connected components.
The components of $K^-(A)$ correspond to those of $G_A$, so \mbox{$\beta_0(K^-(A))=c(G_A)$}.
Lemma~\ref{lem:crux} applies verbatim to the lean complex, because it only concerns triangles, and gives $\beta_2(K^-(A))=k-\rk A$.
With~\eqref{eq:euler} we get
\begin{equation}\label{eq:betti1-KA}
  \beta_1(K^-(A))\;=\;(r-\rk A)\;+\;\bigl(\nnz(A)-r-k+c(G_A)\bigr).
\end{equation}
The second summand is the number of independent cycles of $G_A$, which is typically of order $n$; Edelsbrunner and Parsa make the same observation for their static complex~\cite[Remark~(2) after Theorem~3.1]{EP14}.
So one may hope for a data structure whose cost depends on the homology of $K$ rather than on its size.
Such a structure might be output-sensitive, with a running time that depends on the size of the answer.
We now rule this out: small homology alone does not help.
The lower bound survives on connected complexes with $\beta_1,\beta_2\le1$, and the question is only whether $H_1$ vanishes.

\begin{theorem}[Deciding whether $H_1$ vanishes]\label{thm:acyclic}
Assume the \OMv\ conjecture, and let $\F$ be a field and $\eps>0$.
No data structure with polynomial preprocessing time maintains a subcomplex $K$ of a $2$-dimensional universe with $n$ simplices under simplex insertions and deletions and decides whether $H_1(K;\F)=0$, with amortized update time $O(n^{1/2-\eps})$ and query time $O(n^{1-\eps})$.
This holds even under the promise that at every query $K$ is connected, $\chi(K)=1$, and $\beta_1(K)=\beta_2(K)\in\{0,1\}$.
\end{theorem}

Under the promise, $H_1(K;\F)=0$ holds exactly when $K$ is $\F$-acyclic, so the same lower bound applies to deciding $\F$-acyclicity.
It also applies to approximation: a data structure that returns a number in $[\beta_1(K),c\,\beta_1(K)]$ for some factor $c\ge1$ returns~$0$ exactly when $\beta_1(K)=0$, so it decides whether $H_1(K;\F)=0$.
Hence no multiplicative approximation of~$\beta_1$, with any factor, can be maintained within the time bounds of Theorem~\ref{thm:acyclic}.

The proof combines three ingredients.
The first removes the cycles of the support graph without touching the rank information.
The second is a bordered matrix whose rank deficiency is a single bit.
The third is a standard isolation step that turns that bit into the Boolean answer of \OuMv.

\paragraph{Coning off the support graph.}
Take $P=[r]\times[k]$ in Definition~\ref{def:gadget}, so that column $j$ has the triangle cycles $s^j_0,\dots,s^j_{r+1}$ and the socket between $s^j_t$ and $s^j_{t+1}$ belongs to row $t$ for $t\in[r]$.
Let $a^j_t$ be the vertex $0$ of $s^j_t$, and let $c_i$ be the vertex $0$ of the row cycle $C_i$.
The \emph{spine} $L(A)\subseteq K^-(A)$ is the graph on these vertices with two kinds of edges.
The first kind consists of the edges $a^j_ta^j_{t+1}$ for $0\le t\le r$ and $j\in[k]$.
Such an edge is the edge $a_0b_0$ of the cylinder of the socket between $s^j_t$ and $s^j_{t+1}$, and it survives the removal of the hole, so it lies in $K^-(A)$ whether the socket is sealed or tubed.
The second kind is an \emph{attachment edge} $a^j_ic_i$ for every nonzero entry $A_{ij}$, namely the edge $h_0c_0$ of the tube of position $(i,j)$.
Let $z$ be a new vertex.
The \emph{coned complex} $\widehat K(A)$ is obtained from $K^-(A)$ by adding $z$, an edge $zv$ for every vertex $v$ of $L(A)$, and a triangle $ze$ for every edge $e$ of $L(A)$.
In the coned complex a flip of $A_{ij}$ costs $14$ simplex updates.
These are the $7$ triangles of Lemma~\ref{lem:locality}, the $6$ edges of the tube that are not edges of its hole $h$ or of $C_i$, and the cone triangle over the attachment edge.
When the tube is inserted, we first delete the seal, then insert the six edges, the six tube triangles, and the cone triangle; a deletion of the tube proceeds in the reverse order, so every intermediate family is a subcomplex.
Unlike in the proof of Theorem~\ref{thm:omv}, these flips change the $1$-skeleton.

\begin{lemma}\label{lem:coning}
For every field $\F$ and every matrix $A\in\{0,1\}^{r\times k}$,
\[
  \beta_0(\widehat K(A))=1,\qquad
  \beta_1(\widehat K(A))=r-\rk A,\qquad
  \beta_2(\widehat K(A))=k-\rk A .
\]
\end{lemma}

\begin{proof}
Write $K=K^-(A)$, $L=L(A)$, and $C=z\ast L$ for the cone, so that $\widehat K(A)=K\cup C$ and $K\cap C=L$.
We first show that the inclusion $\iota\colon L\to K$ is injective on $H_1$ and bijective on $\widetilde H_0$.
Define a vertex map $\varphi$ from $K$ to $G_A$ that sends every vertex of $C_i$ to $\rho_i$ and every other vertex of column~$j$ to $\gamma_j$; the latter include the vertices of the tether $\Lambda_j$ and the apices of the two cones of column~$j$.
Every simplex of $K$ lies in one column surface together with its tether filling, or in one row cycle, or in the tube of a position $(i,j)$ with $A_{ij}=1$.
In each case its image is a vertex or an edge of $G_A$, so $\varphi$ is simplicial.
On $L$, the map $\varphi$ collapses the path $a^j_0a^j_1\cdots a^j_{r+1}$ of each column, which is a tree, to the node~$\gamma_j$, and it maps each attachment edge onto the corresponding edge of $G_A$.
Collapsing disjoint trees is a homotopy equivalence, so $(\varphi\circ\iota)_*$ is an isomorphism on all homology groups.
Hence $\iota_*\colon H_1(L)\to H_1(K)$ is injective, with image of dimension $\nnz(A)-r-k+c(G_A)$.
For the same reason $\iota_*\colon\widetilde H_0(L)\to\widetilde H_0(K)$ is injective, and it is an isomorphism because both groups have dimension $c(G_A)-1$: the graph $L$ has $c(G_A)$ components, being homotopy equivalent to $G_A$, and so does $K$.

The cone $C$ is contractible and $H_2(L)=0$, so the Mayer--Vietoris sequence of $K\cup C$ reads
\begin{align*}
  0\to H_2(K)\to H_2(K\cup C)\to H_1(L)&\xrightarrow{\;\iota_*\;}H_1(K)\to H_1(K\cup C)\\
  &\to\widetilde H_0(L)\xrightarrow{\;\iota_*\;}\widetilde H_0(K)\to\widetilde H_0(K\cup C)\to0 .
\end{align*}
Since $\iota_*$ is injective on $H_1$ and bijective on $\widetilde H_0$, we get $H_2(K\cup C)\cong H_2(K)$, $\widetilde H_0(K\cup C)=0$, and $H_1(K\cup C)\cong H_1(K)/\iota_*H_1(L)$.
Substituting $\beta_2(K)=k-\rk A$ and~\eqref{eq:betti1-KA} proves the lemma.
\end{proof}

\paragraph{A bordered matrix with a one-bit rank deficiency.}
We use the Schur-complement identity of Lemma~\ref{lem:bordered}.

For $M\in\{0,1\}^{\ell\times\ell}$ and $u,v\in\{0,1\}^{\ell}$ we use the $(2\ell+1)\times(2\ell+1)$ matrix $\mathrm{Bord}(M;u,v)$ below, whose upper left $2\ell\times2\ell$ block is the invertible matrix~$Y$ to which we apply Lemma~\ref{lem:bordered}:
\[
  \mathrm{Bord}(M;u,v)\;=\;\begin{pmatrix}I&M&0\\0&I&v\\u^{\top}&0&0\end{pmatrix},
  \qquad
  Y=\begin{pmatrix}I&M\\0&I\end{pmatrix},\quad
  Y^{-1}=\begin{pmatrix}I&-M\\0&I\end{pmatrix}.
\]
To keep the field product apart from the Boolean product $u^{\top}Mv$ of \OuMv, we write $\mu_{\F}(u,v)=\sum_{i,j}u_iM_{ij}v_j$, evaluated in~$\F$.
With $x=(u,0)$ and $w=(0,v)$ we get $x^{\top}Y^{-1}w=-\mu_{\F}(u,v)$.
By Lemma~\ref{lem:bordered}, the rank of $\mathrm{Bord}(M;u,v)$ is $2\ell+1$ if $\mu_{\F}(u,v)\ne0$ and $2\ell$ otherwise.
Lemma~\ref{lem:coning} therefore gives
\begin{equation}\label{eq:acyclic-bit}
  \beta_0=1,\qquad \beta_1=\beta_2=\bigl[\mu_{\F}(u,v)=0\bigr]
  \qquad\text{for } K=\widehat K\bigl(\mathrm{Bord}(M;u,v)\bigr).
\end{equation}
So $K$ is $\F$-acyclic exactly when $\mu_{\F}(u,v)\neq0$.

\paragraph{From the field product to the Boolean product.}
The quantity $\mu_{\F}(u,v)$ in~\eqref{eq:acyclic-bit} is computed in $\F$, whereas \OuMv\ asks whether the submatrix $S=M[\supp u,\supp v]$ is nonzero.
These can differ; over $\Ftwo$, for instance, the entries of $S$ may cancel in pairs.
Random restrictions repair this.
\begin{lemma}\label{lem:isolation}
Let $S\in\{0,1\}^{a\times b}$ be nonzero, and let $\xi\in\{0,1\}^a$ and $\eta\in\{0,1\}^b$ be independent and uniformly random.
Then $\Pr[\xi^{\top}S\eta\ne0\text{ in }\F]\ge1/4$ for every field~$\F$.
\end{lemma}
\begin{proof}
The polynomial $f(\xi,\eta)=\sum_{i,j}S_{ij}\xi_i\eta_j$ is multilinear of degree~$2$, and it is nonzero because its coefficients are the entries of~$S$.
A nonzero multilinear polynomial of degree~$d$ does not vanish on at least a $2^{-d}$ fraction of the Boolean cube.
To see this, fix a monomial of maximal degree; for every assignment to the other variables, the restriction to the variables of that monomial is a nonzero multilinear polynomial in $d$ variables, which has a nonzero point among the $2^d$ assignments.
\end{proof}

\begin{proof}[Proof of Theorem~\ref{thm:acyclic}]
Let $M\in\{0,1\}^{\ell\times\ell}$ be an \OuMv\ instance.
The universe $\cU$ is the union of the coned complexes $\widehat K(\mathrm{Bord}(M;u,v))$ over all $u,v\in\{0,1\}^{\ell}$.
It contains the seal and the tube at every border position and has $n=\Theta(\ell^2)$ simplices.
We insert the simplices of $\widehat K(\mathrm{Bord}(M;0,0))$, which costs $O(\ell^2)$ updates.

For a request $(u,v)$, we perform $T=\lceil 8\ln\ell\rceil+1$ independent probes.
A probe draws uniform $\xi\in\{0,1\}^{\supp u}$ and $\eta\in\{0,1\}^{\supp v}$ and sets the border entries accordingly.
It changes at most $2\ell$ entries, each by a constant number of simplex updates.
It then asks whether $H_1(K)=0$.
We answer~$1$ if and only if some probe finds $H_1(K)=0$.

By~\eqref{eq:acyclic-bit}, a probe finds $H_1(K)=0$ exactly when $\xi^{\top}S\eta\ne0$, where $S=M[\supp u,\supp v]$.
If $S=0$, no probe succeeds and the answer~$0$ is correct.
If $S\ne0$, each probe succeeds with probability at least $1/4$ by Lemma~\ref{lem:isolation}.
All $T$ probes then fail with probability at most $(3/4)^T<\ell^{-2}$.
A union bound over the $\ell$ requests bounds the total error probability by $1/\ell$.

At every query, $K$ is connected, $\beta_1(K)=\beta_2(K)\in\{0,1\}$, and $\chi(K)=\beta_0-\beta_1+\beta_2=1$, as promised.
The $\ell$ requests use $O(\ell^2\log\ell)$ updates and $O(\ell\log\ell)$ queries.
A data structure with the stated bounds would therefore solve \OuMv\ with polynomial preprocessing and computation time
\[
  O\bigl(\ell^2\log\ell\cdot n^{1/2-\eps}+\ell\log\ell\cdot n^{1-\eps}+\ell^2\log\ell\bigr)
  \;=\;O\bigl(\ell^{3-\delta}\bigr)\qquad\text{for every }\delta<\min\{2\eps,1\},
\]
where the last term accounts for sampling the probes.
This contradicts the \OMv\ conjecture~\cite{HKNS15}, which covers randomized algorithms with error probability at most $1/3$.
The probes do not depend on the answers of the data structure, so the same argument applies to a randomized data structure that answers all queries of every fixed operation sequence correctly with probability at least $3/4$: the total error probability is then at most $1/4+1/\ell\le1/3$ for $\ell\ge12$.
\end{proof}

\begin{remark}
Over $\Ftwo$, setting the border entries to $u$ and $v$ themselves makes a single query compute the parity $\mu_{\Ftwo}(u,v)$.
Under the analogous conjecture for matrix--vector products over $\Ftwo$, the reduction is deterministic and needs one query per request.
The cone vertex $z$ has degree $\Theta(n)$.
We do not know whether bounded vertex degree, as in Theorem~\ref{thm:bounded-degree}, can be combined with bounded homology.
\end{remark}

\subsection{Scope of the lower bounds}\label{sec:lower-scope}

The bounds do not cover the following settings.
First, in the vineyard model~\cite{CEMS06}, an update transposes two simplices that are consecutive in a total order compatible with the filtration.
A value update of Problem~\ref{prob:sublevel} can require up to $n-1$ such transpositions, so our bounds say nothing about the time per transposition.
Second, for fully dynamic data structures with worst-case bounds, the only lower bounds that we obtain from the \OMv\ conjecture itself are those implied by Theorem~\ref{thm:omv}.
Stronger worst-case bounds rest on the \uMv-hinted \uMv\ conjecture (Corollary~\ref{cor:hinted}), and their exponent depends on that of rectangular matrix multiplication.
Third, the variants of the lower bound are proved in the degrees stated: bounded degree for every $p\ge1$, linear embeddings in $\R^4$ for $\beta_1$ and $\beta_2$, clique and Vietoris--Rips complexes for every $\beta_p$ with $p\ge1$, and the acyclicity test for~$H_1$.
Finally, all these bounds are conditional; the only unconditional lower bound that we know is the $\Omega(\log n)$ cell-probe bound inherited from dynamic connectivity (Section~\ref{sec:easy-conn}).
